\SourcePage{309}
\SGAsection{3}{上同调固有性与泛点附近的局部无环性}
\begin{SGAresult}{定理}{3.1}
设 $S$ 是\Term{generic-point}{泛点}为 $s$ 的\Term{irreducible}{不可约}\Term{scheme}{概形}，$X,Y$ 是两个\Term{finite-presentation}{有限表示}的 $S$-\Term{scheme}{概形}，$f:X\to Y$ 是 $S$-\Term{morphism}{态射}。对每个 $S$-\Term{scheme}{概形} $S'$，以 $Y',X'$ 等记 $Y,X$ 等沿 $S'\to S$ 的\Term{inverse-image}{逆像}。有下列性质：

1) a) 存在 $S$ 的一个非空开子集 $S'$，使得对 $X'$ 上每个\Term{finite}{有限}常值\Term{set-sheaf}{集合层} $F'$，$f'_*F'$ 都是\Term{constructible}{可构造}的，且 $(F',f')$ 相对于 $S'$ 在维数 $\leq0$ 内是\Term{cohomologically-proper}{上同调固有}的。

b) 设 $F$ 是 $X$ 上的\Term{constructible}{可构造}\Term{set-sheaf}{集合层}。则存在 $S$ 的一个非空开子集 $S'$（依赖于 $F$），使得 $f'_*F'$ \Term{constructible}{可构造}，且 $(F',f')$ 相对于 $S'$ 在维数 $\leq0$ 内是\Term{cohomologically-proper}{上同调固有}的。

2) 设 $k(s)$ 的一个\Term{algebraic-closure}{代数闭包} $\bar k$ 上，所有维数 $\leq\dim X_s$ 的\Term{finite-type}{有限型}\Term{scheme}{概形}都可强消解奇点（SGA 5 I 3.1.5）。则还有下列性质：

a) 存在 $S$ 的一个非空开子集 $S'$，使得对 $X'$ 上每个阶与 $S$ 的各\Term{residue-characteristic}{剩余特征}互素的\Term{finite}{有限}常值\Term{group-sheaf}{群层} $F'$，若 $\Phi'$ 是 $F'$-\Term{torsor}{挠子}所成的\Term{stack}{叠}（stack），则 $f'_*\Phi'$ 是 $1$-\Term{constructible}{可构造}的，且 $(F',f')$ 相对于 $S'$ 在维数 $\leq1$ 内是\Term{cohomologically-proper}{上同调固有}的。

b) 设 $L$ 为所有不同于 $S$ 的\Term{residue-characteristic}{剩余特征}的素数组成的集合，$\Phi$ 是 $X$ 上一个 $1$-\Term{constructible}{可构造}的 ind-$L$-\Term{stack}{叠}（0），并且对 $X$ 上每个\Term{etale}{平展}\Term{scheme}{概形} $X_1$ 以及 $\Phi_{X_1}$ 中的任意两个对象 $x,x_1$，层 $\underline{\operatorname{Hom}}_{X_1}(x,x_1)$ 都\Term{constructible}{可构造}。再设 $S$ \Term{locally-noetherian}{局部 Noether}。则存在 $S$ 的一个非空开子集 $S'$，使得 $f'_*\Phi'$ 是 $1$-\Term{constructible}{可构造}的，$(f'_*\Phi')_{Y_1}$ 的任一\Term{fiber}{纤维}中的任意两个对象 $y,y_1$ 的 $\underline{\operatorname{Hom}}_{Y_1}(y,y_1)$ 都\Term{constructible}{可构造}，而且 $(\Phi',f')$ 相对于 $S'$ 在维数 $\leq1$ 内是\Term{cohomologically-proper}{上同调固有}的。
\end{SGAresult}
\begin{proof}
可设 $S$ \Term{affine}{仿射}；由 SGA 4 VIII 1.1，可设 $S$ \Term{integral}{整}；最后，利用取极限，可设 $S$ 是 $\mathbf Z$ 上某个\Term{finite-type}{有限型}代数的谱，因而 $S$ 是 Noether \Term{scheme}{概形}。问题在 $Y$ 上具有局部性，因此可设 $Y$ \Term{affine}{仿射}。此外，只需在\Term{finite}{有限}\Term{group-extension}{扩张} $S'\to S$ 后证明定理，其中 $S'$ 是\Term{integral-scheme}{整概形}，而 $S'\to S$ 是\Term{etale-morphism}{平展态射}与\Term{finite}{有限}、\Term{radicial}{根性}、\Term{surjectivity}{满射}\Term{morphism}{态射}的复合。

\textbf{1) 常值\Term{set-sheaf}{集合层}的情形。}

\textbf{1) 1. 归约到 $X$ 在 $S$ 上\Term{normal}{正规}的情形。}设 $X_{1\bar s}$ 为 $(X_{\bar s})_{\mathrm{red}}$ 的\Term{normalization}{正规化}。将 $S$ 缩小为一个非空开子集并对 $S$ 作\Term{radicial}{根性}\Term{group-extension}{扩张}后，可设 $X_{1\bar s}$ 来自 $S$ 上的\Term{normal}{正规}\Term{scheme}{概形} $X_1$，而\Term{morphism}{态射} $X_{1\bar s}\to X_{\bar s}$ 来自\Term{finite}{有限}\Term{surjectivity}{满射}
\SourcePage{310}
$p:X_1\to X$（EGA IV 8.8.2、9.6.1）。设定理已对 $fp$ 成立。将 $S$ 缩小为开子集后，可设对 $X$ 上每个常值\Term{set-sheaf}{集合层} $F$，$(p^*F,fp)$ 相对于 $S$ 在维数 $\leq0$ 内\Term{cohomologically-proper}{上同调固有}，且 $f_*p_*(p^*F)$ \Term{constructible}{可构造}。由 1.9，$(p_*p^*F,f)$ 相对于 $S$ 在维数 $\leq0$ 内也\Term{cohomologically-proper}{上同调固有}。\Term{morphism}{态射}
\[
 F\longrightarrow p_*p^*F=G
\]
是\Term{monomorphism}{单态射}。由于 $f_*F$ 是 $f_*G$ 的\Term{subsheaf}{子层}，已可推出 $f_*F$ \Term{constructible}{可构造}（SGA 4 IX 2.9(ii)），且 $(F,f)$ 相对于 $S$ 在维数 $\leq-1$ 内\Term{cohomologically-proper}{上同调固有}。

设 $X_2=X_1\times_X X_1$，$q:X_2\to X$ 是典范\Term{morphism}{态射}。由 1.11 1)，有\Term{exact-sequence}{正合列}
\[
\begin{tikzcd}
F\arrow[r]&G\arrow[r,shift left=.45ex]\arrow[r,shift right=.45ex]&q_*q^*F.
\end{tikzcd}
\]
将刚才所得结论应用于 $fq$，可知将 $S$ 缩小为非空开子集后，对 $X$ 上每个常值\Term{set-sheaf}{集合层} $F$，$(q^*F,fq)$ 相对于 $S$ 在维数 $\leq-1$ 内\Term{cohomologically-proper}{上同调固有}，从而 $(q_*q^*F,f)$ 也如此。于是由 1.13 1)，$(F,f)$ 相对于 $S$ 在维数 $\leq0$ 内\Term{cohomologically-proper}{上同调固有}。

\textbf{1) 2. 归约到 $X$ 在 $S$ 上\Term{normal}{正规}且\Term{affine}{仿射}的情形。}

设 $U_s$ 是 $X_s$ 中的稠密\Term{affine}{仿射}开子集。将 $S$ 缩小为非空开子集后，可设 $U_s\to X_s$ 提升为相对于 $S$ 在\Term{scheme}{概形}意义下\Term{dominant}{占优}的\Term{open-immersion}{开浸入} $i:U\to X$（EGA IV 8.9.1）。由于 $X\to S$ 是\Term{normal}{正规}\Term{morphism}{态射}，由 SGA 2 XIV 1.18，
\[
\operatorname{prof\text{ét}}_{S-U}(X)\geq2;
\]
因而对 $X$ 上每个\Term{constant-sheaf}{常值层} $F$，典范\Term{morphism}{态射}
\[F\longrightarrow i_*i^*F\]
是\Term{isomorphism}{同构}。因此，若定理已对 $fi$ 和 $i$ 成立，则将 $S$ 缩小为非空开子集后，$(i^*F,i)$ 与 $(i^*F,fi)$ 都相对于 $S$ 在维数 $\leq0$ 内\Term{cohomologically-proper}{上同调固有}。由 1.6 2)，$(F,f)$ 也如此。再由 $f_*F=(fi)_*(i^*F)$ \Term{constructible}{可构造}，归约完成。

\textbf{1) 3. 证明的结束。}

可设 $S$ \Term{normal}{正规}（EGA IV 7.8.3）。可取 $X_s$ 的一个\Term{compactification}{紧化}：
\SourcePage{311}
\[
\begin{tikzcd}
X_s\arrow[r,"j_s"]\arrow[d,"f_s"']&P_s\arrow[dl,"g_s"]\\
Y_s&
\end{tikzcd}
\]
其中 $j_s$ 为\Term{dominant}{占优}\Term{open-immersion}{开浸入}，$g_s$ 为\Term{proper}{固有}\Term{morphism}{态射}。对 $k(s)$ 作\Term{radicial}{根性}\Term{group-extension}{扩张}，再以 $P_s$ 的\Term{normalization}{正规化}替换它；这不改变 $X_s$，并可使 $P_s$ 几何\Term{normal}{正规}。将 $S$ 缩小为非空开子集并作\Term{radicial}{根性}\Term{surjectivity}{满射}\Term{group-extension}{扩张}后，可设上图来自一个图表
\[
\begin{tikzcd}
X\arrow[r,"j"]\arrow[d,"f"']&P\arrow[dl,"g"]\\
Y&
\end{tikzcd}
\]
其中 $P$ 在 $S$ 上\Term{normal}{正规}，$j$ 为相对于 $S$ 在\Term{scheme}{概形}意义下\Term{dominant}{占优}的\Term{open-immersion}{开浸入}，$g$ 为\Term{proper-morphism}{固有态射}（EGA IV 6.9.1、9.9.4、9.6.1）。对 $X$ 上每个值为 $C$ 的\Term{finite}{有限}常值\Term{set-sheaf}{集合层} $F$，$j_*F$ 是值为 $C$ 的\Term{constant-sheaf}{常值层}（SGA 4 2.14.1），并且在任意\Term{base-change}{基变换} $S'\to S$ 后仍然如此。故 $(F,j)$ 相对于 $S$ 在维数 $\leq0$ 内\Term{cohomologically-proper}{上同调固有}。由于 $g$ \Term{proper}{固有}，$(F,f)$ 亦如此（1.8）。又由于 $g$ \Term{proper}{固有}，$f_*F=g_*C_P$ \Term{constructible}{可构造}，1) a) 得证。

\textbf{2) \Term{constructible}{可构造}\Term{set-sheaf}{集合层}的情形。}

设 $F$ 是 $X$ 上的\Term{constructible}{可构造}\Term{set-sheaf}{集合层}。由 SGA 4 IX 2.14(ii)，存在\Term{finite}{有限}族\Term{morphism}{态射} $p_i:Z_i\to X$ 以及每个 $Z_i$ 上的\Term{finite}{有限}常值\Term{set-sheaf}{集合层} $C_i$，使得有\Term{monomorphism}{单态射}
\[j:F\longrightarrow\prod_i p_{i*}C_i=G.\]
由 1) a)，缩小 $S$ 为非空开子集后，可设 $(C_i,fp_i)$ 相对于 $S$ 在维数 $\leq0$ 内都\Term{cohomologically-proper}{上同调固有}，且 $f_*p_{i*}C_i$ \Term{constructible}{可构造}。由 1.9，可先推出 $(G,f)$ 相对于 $S$ 在维数 $\leq0$ 内\Term{cohomologically-proper}{上同调固有}；因而 $(F,f)$ 在维数 $\leq-1$ 内\Term{cohomologically-proper}{上同调固有}，且 $f_*F$ \Term{constructible}{可构造}。令 $K=G\amalg_F G$ 为\Term{amalgamated-sum}{合并和}。由于 $F,G$ \Term{constructible}{可构造}，$K$ 也\Term{constructible}{可构造}。由前面的结论，缩小 $S$ 为非空开子集后，可设 $(K,f)$ 相对于 $S$ 在维数 $\leq-1$ 内\Term{cohomologically-proper}{上同调固有}。再由 1.13 1)，$(F,f)$ 相对于 $S$ 在维数 $\leq0$ 内\Term{cohomologically-proper}{上同调固有}。

\SourcePage{312}
\textbf{3) 常值\Term{group-sheaf}{群层}的情形。}

若 $F$ 是 $X$ 上的常值\Term{group-sheaf}{群层}，以 $\Phi$ 记 $F$-\Term{torsor}{挠子}所成的\Term{stack}{叠}。

\textbf{3) 1.} 先证明：缩小 $S$ 为非空开子集后，对 $X$ 上每个阶与 $S$ 的\Term{residue-characteristic}{剩余特征}互素的常值\Term{group-sheaf}{群层} $F$，$(F,f)$ 相对于 $S$ 在维数 $\leq0$ 内\Term{cohomologically-proper}{上同调固有}，且 $f_*\Phi$ \Term{constructible}{可构造}。

为此先归约到 $X$ 在 $S$ 上\Term{smooth}{光滑}的情形。对 $k(s)$ 作\Term{finite}{有限}\Term{group-extension}{扩张}后——这是允许的，因为它可视为\Term{etale}{平展}\Term{group-extension}{扩张}与\Term{radicial}{根性}\Term{group-extension}{扩张}的复合——可取\Term{proper}{固有}\Term{surjectivity}{满射} $p_s:X_{1s}\to X_s$，其中 $X_{1s}$ 是 $S$ 上与 $X_s$ 同维的\Term{smooth}{光滑}\Term{scheme}{概形}。再缩小 $S$ 为非空开子集，可设 $p_s$ 来自\Term{proper}{固有}\Term{surjectivity}{满射} $p:X_1\to X$，其中 $X_1$ 在 $S$ 上\Term{smooth}{光滑}（EGA IV 9.6.1、12.1.6）。设 $X_2=X_1\times_X X_1$，$q:X_2\to X$ 是典范\Term{morphism}{态射}。由 1.11 2)，有 $X$ 上的叠的\Term{exact}{正合}图表
\[
\begin{tikzcd}
\Phi\arrow[r]&p_*p^*\Phi\arrow[r,shift left=.45ex]\arrow[r,shift right=.45ex]&q_*q^*\Phi.
\end{tikzcd}
\]
设定理已在\Term{smooth}{光滑}情形下成立，则先可设 $f_*p_*p^*\Phi$ \Term{constructible}{可构造}，由下文 3.1.1，$f_*\Phi$ 亦\Term{constructible}{可构造}。此外，由 1.6 2)，可设 $(p_*p^*\Phi,f)$ 相对于 $S$ 在维数 $\leq0$ 内\Term{cohomologically-proper}{上同调固有}，故 $(\Phi,f)$ 在维数 $\leq-1$ 内\Term{cohomologically-proper}{上同调固有}。于是又可设 $(q_*q^*\Phi,f)$ 相对于 $S$ 在维数 $\leq-1$ 内\Term{cohomologically-proper}{上同调固有}；由 1.12 1)，这就蕴含 $(\Phi,f)$ 在维数 $\leq0$ 内\Term{cohomologically-proper}{上同调固有}。

再像 1) 2 那样，归约到 $X$ 在 $S$ 上\Term{smooth}{光滑}且\Term{affine}{仿射}的情形。取 $X$ 的\Term{compactification}{紧化}
\[
\begin{tikzcd}
X\arrow[r,"i"]\arrow[d,"f"']&P\arrow[dl,"q"]\\
Y&
\end{tikzcd}
\]
其中 $i$ 为\Term{dominant}{占优}\Term{open-immersion}{开浸入}，$q$ 为\Term{proper-morphism}{固有态射}。因为 $\dim P_s=\dim X_s$，可对 $P_s$ 应用消解奇点的假设。对 $S$ 作\Term{etale}{平展}\Term{group-extension}{扩张}与\Term{radicial}{根性}\Term{group-extension}{扩张}后，可取\Term{proper-morphism}{固有态射} $r:Z\to P$，其中 $Z$ 在 $S$ 上\Term{smooth}{光滑}，$r^{-1}(X)\simeq X$，且 $r^{-1}(X)$ 是 $Z$ 中一个相对于 $S$ 具有\Term{normal-crossings}{正规交叉}的\Term{divisor}{除子}的补集。于是每个 $F$-\Term{torsor}{挠子}都在 $Z$ 上相对于 $S$ \Term{tamely-ramified}{温和分歧}（2.3 b)）。由 2.7，$(F,f)$ 相对于 $S$ 在维数 $\leq0$ 内\Term{cohomologically-proper}{上同调固有}，断言得证。

\SourcePage{313}
\textbf{3) 2. 归约到 $X$ 在 $S$ 上\Term{smooth}{光滑}的情形。}

对 $k(s)$ 作\Term{finite}{有限}\Term{group-extension}{扩张}后，可取\Term{proper}{固有}\Term{surjectivity}{满射} $p_s:X_{1s}\to X$，其中 $X_{1s}$ 是 $s$ 上的\Term{smooth}{光滑}\Term{scheme}{概形}，并可设 $p_s$ 来自\Term{proper}{固有}\Term{surjectivity}{满射} $p:X_1\to X$，其中 $X_1$ 在 $S$ 上\Term{smooth}{光滑}。设定理已对 $fp$ 成立，来证明它对 $f$ 成立。设 $F$ 为 $X$ 上阶与 $S$ 的\Term{residue-characteristic}{剩余特征}互素的\Term{finite}{有限}常值\Term{group-sheaf}{群层}，$\Phi$ 为 $F$-\Term{torsor}{挠子}所成的叠。置
\[X_2=X_1\times_X X_1,\qquad X_3=X_1\times_X X_1\times_X X_1,\]
并以 $q:X_2\to X$、$r:X_3\to X$ 记典范\Term{morphism}{态射}。由 1.11 2)，有叠的\Term{exact}{正合}图表
\[
\begin{tikzcd}[column sep=small]
\Phi\arrow[r]&p_*p^*\Phi\arrow[r,shift left=.45ex]\arrow[r,shift right=.45ex]&q_*q^*\Phi\arrow[r,shift left=.8ex]\arrow[r]\arrow[r,shift right=.8ex]&r_*r^*\Phi.
\end{tikzcd}
\]
要证明 $(\Phi,f)$ 相对于 $S$ 在维数 $\leq1$ 内\Term{cohomologically-proper}{上同调固有}，只需证明 $(p_*p^*\Phi,f)$ 亦如此，$(q_*q^*\Phi,f)$ 相对于 $S$ 在维数 $\leq0$ 内\Term{cohomologically-proper}{上同调固有}，且 $(r_*r^*\Phi,f)$ 在维数 $\leq-1$ 内\Term{cohomologically-proper}{上同调固有}（1.12 2)）。由上面的 3) 1，可设对每个\Term{finite}{有限}常值\Term{group-sheaf}{群层} $F$，$(q^*\Phi,fq)$、$(q^*\Phi,q)$、$(r^*\Phi,fr)$、$(r^*\Phi,r)$ 都相对于 $S$ 在维数 $\leq0$ 内\Term{cohomologically-proper}{上同调固有}。由 1.6 2)，$(q_*q^*\Phi,f)$、$(r_*r^*\Phi,f)$ 在维数 $\leq0$ 内\Term{cohomologically-proper}{上同调固有}。又设定理已在\Term{smooth}{光滑}情形下成立，则 $(p^*\Phi,fp)$、$(p^*\Phi,p)$，进而 $(p_*p^*\Phi,f)$ 都相对于 $S$ 在维数 $\leq1$ 内\Term{cohomologically-proper}{上同调固有}（1.6 2)）。故 $(\Phi,f)$ 相对于 $S$ 在维数 $\leq1$ 内\Term{cohomologically-proper}{上同调固有}。

此外，按假设 $f_*p_*p^*\Phi$ 是 $1$-\Term{constructible}{可构造}的；由 3.1，可设 $f_*q_*q^*\Phi$ \Term{constructible}{可构造}。由下文 3.1.1，$f_*\Phi$ 是 $1$-\Term{constructible}{可构造}的。

\textbf{3) 3. 归约到 $X$ 在 $S$ 上\Term{smooth}{光滑}且\Term{affine}{仿射}的情形。}

由 3) 2，可设 $X$ 在 $S$ 上\Term{smooth}{光滑}。设 $(U_i)_{i\in I}$ 为 $X$ 的一个\Term{finite}{有限}\Term{affine}{仿射}开\Term{covering}{覆盖}，$X_1$ 为 $U_i$ 的\Term{direct-sum}{直和}，$p:X_1\to X$ 为典范\Term{morphism}{态射}。由于 $p$ 是变动\Term{scheme}{概形}上\Term{finite-type}{有限型}\Term{etale}{平展}层所成\Term{category}{范畴}的\Term{effective-descent}{有效下降}\Term{morphism}{态射}，像 3) 2 那样可知：若定理已对 $U_i$，即对 $X_1$ 成立，则对 $X$ 也成立。

\textbf{3) 4. $X$ 在 $S$ 上\Term{smooth}{光滑}且\Term{affine}{仿射}的情形。}

像 3) 1 那样可知，将 $S$ 缩小为非空开子集，并作\Term{etale}{平展}\Term{surjectivity}{满射}\Term{group-extension}{扩张}和\Term{radicial}{根性}\Term{surjectivity}{满射}\Term{group-extension}{扩张}后，可取\Term{commutative-diagram}{交换图}
\[
\begin{tikzcd}
X\arrow[r,"i"]\arrow[d,"f"']&P\arrow[dl,"q"]\\
Y&
\end{tikzcd}
\]
其中 $P$ 在 $S$ 上\Term{smooth}{光滑}，$X$ 是 $P$ 中一个相对于 $S$ 具有\Term{normal-crossings}{正规交叉}的\Term{divisor}{除子}的补集，$g$ 为\Term{proper-morphism}{固有态射}。若 $F$ 是阶与 $S$ 的\Term{residue-characteristic}{剩余特征}互素的\Term{constant-sheaf}{常值层}，则每个 $F$-\Term{torsor}{挠子}都在 $P$ 上相对于 $S$ \Term{tame-ramification}{温和分歧}（2.3 b)）。
\SourcePage{314}
由 2.7，$(F,f)$ 相对于 $S$ 在维数 $\leq1$ 内\Term{cohomologically-proper}{上同调固有}，且 $f_*\Phi$ \Term{constructible}{可构造}。

\textbf{4) 2) b) 的证明。}

\textbf{4) 1. $\Phi$ 为\Term{gerbe}{凝集}的情形。}

可取\Term{finite-type}{有限型}\Term{etale}{平展}\Term{surjectivity}{满射} $p:X_1\to X$，使 $p^*\Phi$ 为\Term{trivial-bundle}{平凡}\Term{gerbe}{凝集}（gerbe）。由\Term{descent}{下降}，像 3) 2 那样，只需对 $X_1$、$X_1\times_X X_1$ 和 $X_1\times_X X_1\times_X X_1$ 证明定理。于是归约到 $\Phi$ 是某个\Term{constructible}{可构造}\Term{group-sheaf}{群层} $F$ 的\Term{torsor}{挠子}所成\Term{gerbe}{凝集}的情形，其中 $F$ 的\Term{fiber}{纤维}的阶与 $S$ 的\Term{residue-characteristic}{剩余特征}互素。

由 SGA 4 IX 2.14，存在\Term{finite}{有限}族\Term{finite-morphism}{有限态射} $p_i:Z_i\to X$，以及每个 $Z_i$ 上阶与 $S$ 的\Term{residue-characteristic}{剩余特征}互素的\Term{finite}{有限}常值\Term{group-sheaf}{群层} $C_i$，使得有\Term{monomorphism}{单态射}
\[j:F\longrightarrow\prod_i p_{i*}C_i=G.\]
设 $\Phi_i$ 是 $C_i$-\Term{torsor}{挠子}所成的叠，$\Psi$ 是 $G$-\Term{torsor}{挠子}所成的叠。由 2) a)，缩小 $S$ 为非空开子集后，可设 $(C_i,fp_i)$ 相对于 $S$ 在维数 $\leq1$ 内都\Term{cohomologically-proper}{上同调固有}，且叠 $f_*p_{i*}\Phi_i$ 都是 $1$-\Term{constructible}{可构造}的。由 1.9，$(p_{i*}C_i,f)$ 相对于 $S$ 在维数 $\leq1$ 内\Term{cohomologically-proper}{上同调固有}；$(G,f)$ 也如此。此外，$p_{i*}\Phi_i$ 等价于群 $p_{i*}C_i$ 的\Term{torsor}{挠子}所成的叠（SGA 4 VIII 5.8），故 $f_*\Psi$ 为 $1$-\Term{constructible}{可构造}。

由于 $R^1f_*G$ \Term{constructible}{可构造}，存在一个由\Term{finite-type}{有限型}\Term{etale}{平展} $Y$-\Term{scheme}{概形} $T$ 表示的层以及一个\Term{epimorphism}{满态射}
\[a:T\longrightarrow R^1f_*G\]
（SGA 4 IX 2.7）；并可设 $T(T)$ 中恒等\Term{section}{截面}的像由 $X\times_Y T=X_T$ 上、群为 $G|_{X_T}$ 的\Term{torsor}{挠子} $Q$ 定义。设 $f_T:X_T\to T$ 为典范\Term{morphism}{态射}，$F_T=F|_{X_T}$，其余记号类似。由 1) b)，缩小 $S$ 为非空开子集后，可设 $(Q/F_T,f_T)$ 相对于 $S$ 在维数 $\leq0$ 内\Term{cohomologically-proper}{上同调固有}，且 $f_{T*}(Q/F_T)$ \Term{constructible}{可构造}。由 3.1.2，$f_*\Psi$ \Term{constructible}{可构造}。

来证明 $(F,f)$ 相对于 $S$ 在维数 $\leq1$ 内\Term{cohomologically-proper}{上同调固有}。由 1.13 2)，只需证明：对 $Y$ 上每个\Term{etale}{平展}\Term{scheme}{概形} $Y_1$ 和 $X_1=X\times_Y Y_1$ 上每个\Term{torsor}{挠子} $Q_1$，若 $f_1:X_1\to Y_1$ 为典范\Term{morphism}{态射}，则 $(Q_1/F_1,f_1)$ 相对于 $S$ 在维数 $\leq0$ 内\Term{cohomologically-proper}{上同调固有}。按 $T$ 的定义，$Q_1$ 在 $Y_1$ 的\Term{etale-topology}{平展拓扑}下局部是 $Q$ 的\Term{inverse-image}{逆像}，归约得证。

\SourcePage{315}
\textbf{4) 2. 一般情形。}

由引理 6.1.1、4) 1 和 1) a)，缩小 $S$ 为非空开子集后，可设 $S(f_*\Phi)$ \Term{constructible}{可构造}，且 $(S\Phi,f)$ 相对于 $S$ 在维数 $\leq0$ 内\Term{cohomologically-proper}{上同调固有}。于是存在一个由\Term{finite-type}{有限型}\Term{etale}{平展} $Y$-\Term{scheme}{概形} $T$ 表示的层以及一个\Term{epimorphism}{满态射}
\[a:T\longrightarrow S(f_*\Phi).\]
沿用 4) 1 的记号，并置 $Z=T\times_Y T$、$X_Z=X\times_Y Z$，以 $f_Z:X_Z\to Z$ 记典范\Term{morphism}{态射}。可选择 $T$，使 $T(T)$ 的恒等\Term{section}{截面}经 $a$ 所得的像 $q$ 由 $(f_*\Phi)_T=\Phi_{X_T}$ 中一个对象 $p$ 定义。设 $p_1,p_2$（分别地 $q_1,q_2$）为 $p$（分别地 $q$）沿 $Z\to T$ 的两个投影的\Term{inverse-image}{逆像}。缩小 $S$ 为非空开子集后，可设 $(\underline{\operatorname{Aut}}_{X_T}(p),f_T)$ 相对于 $S$ 在维数 $\leq1$ 内\Term{cohomologically-proper}{上同调固有}，$(\underline{\operatorname{Hom}}_{X_Z}(p_1,p_2),f_Z)$ 相对于 $S$ 在维数 $\leq0$ 内\Term{cohomologically-proper}{上同调固有}，且
\[
f_{T*}(\underline{\operatorname{Aut}}_{X_T}(p)),\qquad
f_{Z*}(\underline{\operatorname{Hom}}_{X_Z}(p_1,p_2))
\]
\Term{constructible}{可构造}（a) 1) 与 4) 1)）。

先由此得到：对 $Y$ 上每个\Term{etale}{平展}\Term{scheme}{概形} $Y_1$ 及 $(f_*\Phi)_{Y_1}$ 中任意两个对象 $y,y_1$，层 $\underline{\operatorname{Aut}}_{Y_1}(y)$（分别地 $\underline{\operatorname{Hom}}_{Y_1}(y,y_1)$）\Term{constructible}{可构造}，因为在 $Y_1$ 的\Term{etale-topology}{平展拓扑}下，它局部是 $f_{T*}(\underline{\operatorname{Aut}}_{X_T}(p))$（分别地 $f_{Z*}(\underline{\operatorname{Hom}}_{X_Z}(p_1,p_2))$）的\Term{inverse-image}{逆像}。

还须证明 $(\Phi,f)$ 相对于 $S$ 在维数 $\leq1$ 内\Term{cohomologically-proper}{上同调固有}。为此只需证明：对每个 $S'$-\Term{scheme}{概形} $S$ 和 $Y'$ 的每个\Term{geometric-point}{几何点} $\bar y'$，若以 $\bar Y$ 记 $Y$ 在 $\bar y'$ 处的\Term{strict-localization}{严格局部化}，$\bar X=X\times_Y\bar Y$，其余记号类似，则典范\Term{functor}{函子}
\[\bar\varphi:\Phi(\bar X)\longrightarrow\Phi'(\overline{X'})\]
是\Term{equivalence-of-categories}{范畴等价}。

来证明 $\bar\varphi$ \Term{fully-faithful}{完全忠实}。设 $x,y\in\Phi(\bar X)$，$x',y'$ 为它们在 $\Phi'(\overline{X'})$ 中的像。要证明典范\Term{morphism}{态射}
\[
\operatorname{Hom}_{\bar X}(x,y)\longrightarrow
\operatorname{Hom}_{\overline{X'}}(x',y')\tag{*}
\]
是\Term{bijection}{双射}。由 $T$ 的定义，存在两个从 $\bar Y$ 到 $T$ 的\Term{morphism}{态射}，使 $x,y$ 分别是 $p$ 沿这两个\Term{morphism}{态射}的\Term{inverse-image}{逆像}。这等价于有\Term{morphism}{态射} $\bar Y\to Z$，进而有\Term{morphism}{态射} $h:\bar X\to X_Z$，使
\[h^*(p_1)=x,\qquad h^*(p_2)=y.\]
故有典范\Term{isomorphism}{同构}
\[\underline{\operatorname{Hom}}_{\bar X}(x,y)=h^*(\underline{\operatorname{Hom}}_{X_Z}(p_1,p_2)).\]
由于 $(\underline{\operatorname{Hom}}_{X_Z}(p_1,p_2),f_Z)$ 相对于 $S$ 在维数 $\leq0$ 内\Term{cohomologically-proper}{上同调固有}，$(\underline{\operatorname{Hom}}_{\bar X}(x,y),\bar f)$ 也如此，故 (*) 为\Term{bijection}{双射}。

\SourcePage{316}
再证明 $\bar\varphi$ \Term{essentially-surjective}{本质满}。设 $x'\in\Phi'(\overline{X'})$。由于 $(S\Phi,f)$ 相对于 $S$ 在维数 $\leq0$ 内\Term{cohomologically-proper}{上同调固有}，典范\Term{morphism}{态射}
\[H^0(\bar X,\overline{S\Phi})\longrightarrow H^0(\overline{X'},\overline{S\Phi'})\]
是\Term{bijection}{双射}。设 $G'$ 为 $x'$ 生成的 $\Phi'$ 的\Term{maximal-subgerbe}{极大子凝集}，则存在 $\Phi$ 的\Term{maximal-subgerbe}{极大子凝集} $G$，其在 $\overline{X'}$ 上的\Term{inverse-image}{逆像}为 $G'$。由 4) 1，$(G,f)$ 相对于 $S$ 在维数 $\leq1$ 内\Term{cohomologically-proper}{上同调固有}；故典范\Term{functor}{函子}
\[G(\bar X)\longrightarrow G'(\overline{X'})\]
为\Term{category-equivalence}{范畴等价}。这证明存在 $\Phi(\bar X)$ 中的一个元素 $x$，其在 $\Phi'(\overline{X'})$ 中的像\Term{isomorphism}{同构}于 $x'$，定理得证。
\end{proof}

\begin{SGAresult}{引理}{3.1.1}
设 $S$ 是\Term{locally-noetherian}{局部 Noether} \Term{scheme}{概形}，
\[
\begin{tikzcd}
\Phi\arrow[r,"p"]&\Phi_1\arrow[r,shift left=.5ex,"p_1"]\arrow[r,shift right=.5ex,"p_2"']&\Phi_2
\end{tikzcd}
\]
是 $S$ 上叠的\Term{exact}{正合}图表（1.10.1）。若 $\Phi_1$ \Term{constructible}{可构造}，则 $\Phi$ 也\Term{constructible}{可构造}。若对 $S$ 上每个\Term{etale}{平展}\Term{scheme}{概形} $S'$ 以及 $(\Phi_1)_{S'}$ 中的任意两个对象 $x_1,y_1$，层 $\underline{\operatorname{Hom}}_{S'}(x_1,y_1)$ 都\Term{constructible}{可构造}，则对 $\Phi_{S'}$ 中任意两个对象 $x,y$，$\underline{\operatorname{Hom}}_{S'}(x,y)$ 也\Term{constructible}{可构造}。若 $\Phi_1$ 是 $1$-\Term{constructible}{可构造}的（0），且 $\Phi_2$ \Term{constructible}{可构造}，则 $\Phi$ 是 $1$-\Term{constructible}{可构造}的。
\end{SGAresult}
\begin{proof}
对 $S$ 上每个\Term{etale}{平展}\Term{scheme}{概形} $S'$ 和 $\Phi_{S'}$ 中每个对象 $x$，有\Term{monomorphism}{单态射}
\[\underline{\operatorname{Aut}}_{S'}(x)\longrightarrow\underline{\operatorname{Aut}}_{S'}(p(x)).\]
由 SGA 4 IX 2.9，若 $\Phi_1$ \Term{constructible}{可构造}，则 $\Phi$ \Term{constructible}{可构造}；引理的第二个断言同样可证。

现设 $\Phi_1$ 为 $1$-\Term{constructible}{可构造}，$\Phi_2$ \Term{constructible}{可构造}。\Term{morphism}{态射} $p$ 在\Term{maximal-subgerbe}{极大子凝集}层上诱导\Term{morphism}{态射}
\[\varphi:S\Phi\longrightarrow S\Phi_1.\]
设 $G$ 为 $S\Phi$ 经 $\varphi$ 的像；由 SGA 4 IX 2.9，$G$ 是\Term{constructible}{可构造}层。于是存在由\Term{finite-type}{有限型}\Term{etale}{平展} $S$-\Term{scheme}{概形} $T$ 表示的层以及\Term{epimorphism}{满态射}
\[a:T\longrightarrow G\]
（SGA 4 IX 2.7）。还可选择 $T$，使 $T(T)$ 中恒等\Term{section}{截面}的像 $y$ 由 $(\Phi_1)_T$ 中形如 $x_1=p(x)$ 的对象定义，其中 $x\in\Phi_T$。

只需证明：对 $S$ 的每个点 $s$，$\overline{\{s\}}$ 中存在非空开子集 $U$，使 $S\Phi|_U$ 为\Term{locally-constant}{局部常值}\Term{constructible}{可构造}层。设 $s\in S$，$\bar s$ 为其上方的\Term{geometric-point}{几何点}，$\bar y_1,\ldots,\bar y_n$ 为 $G_{\bar s}$ 的各元素。由 $T$ 的定义，存在\Term{morphism}{态射} $h_i:\bar s\to T$，使 $h_i^*(y)=\bar y_i$。设 $S'$ 为 $n$ 个\Term{isomorphism}{同构}于 $T$ 的\Term{scheme}{概形}在 $S$ 上的\Term{fiber-product}{纤维积}，$h:\bar s\to S'$ 为 $h_i$ 的\Term{fiber-product}{纤维积}，$y_i$（分别地 $x_i$）为 $y$（分别地 $x$）沿 $S'\to T$ 的第 $i$ 个投影的\Term{inverse-image}{逆像}。
\SourcePage{317}
设 $F_i$ 为 $S\Phi|_{S'}$ 中 $y_i$ 的\Term{inverse-image}{逆像}\Term{subsheaf}{子层}，来证明 $F_i$ \Term{constructible}{可构造}。层 $F_i$ 是层 $F'_i$ 的商，其中，对 $S'$ 上每个\Term{etale}{平展}\Term{scheme}{概形} $S''$，
\[
\begin{aligned}
F'_i(S'')=\{&\Phi_{S''}\text{ 中配备同构 }i:p(z)\simeq p(x_i|_{S''})\\
&\text{的对象 }z\text{ 的同构类}\}.
\end{aligned}
\]
只需证明 $F'_i$ \Term{constructible}{可构造}。置 $z_i=p_1p(x_i)$、$z'_i=p_2p(x_i)$，则有\Term{monomorphism}{单态射}
\[\Psi_i:F'_i\longrightarrow\underline{\operatorname{Isom}}_{S'}(z_i,z'_i).\]
它将 $S'$ 上每个\Term{etale}{平展}\Term{scheme}{概形} $S''$ 和 $\Phi_{S''}$ 中配备\Term{isomorphism}{同构} $i:p(z)\simeq p(x_i|_{S''})$ 的对象 $z$，对应到由下图交换性所确定的 $z_i\to z'_i$ 的\Term{isomorphism}{同构}：
\[
\begin{tikzcd}
p_1p(z)\arrow[r,"p_1(i)"]\arrow[d,"j"']&z_i\arrow[d]\\
p_2p(z)\arrow[r,"p_2(i)"']&z'_i.
\end{tikzcd}
\]
这里 $j$ 为与\Term{exact}{正合}图表相伴的典范\Term{morphism}{态射}。

\Term{morphism}{态射} $\Psi_i$ 为\Term{injectivity}{单射}，因为对配备\Term{isomorphism}{同构} $i:p(z)\xrightarrow{\sim}p(x_i|_{S''})$、$i':p(z')\xrightarrow{\sim}p(x_i|_{S''})$ 的两个对象 $z,z'$，它们在 $\operatorname{Isom}_{S''}(z_i,z'_i)$ 中定义相同元素，当且仅当 $p_1(i'^{-1}i)=p_2(i'^{-1}i)$，亦即 $i'^{-1}i$ 来自一个\Term{isomorphism}{同构} $z\to z'$。层 $F'_i$ 是 $\underline{\operatorname{Isom}}_{S'}(z_i,z'_i)$ 的\Term{subsheaf}{子层}，故\Term{constructible}{可构造}。

可取 $\overline{\{s\}}$ 中一个非空开子集 $U$，使 $y_1|_U,\ldots,y_n|_U$ 生成 $G$；由下文引理 6.1.2，$S\Phi|_U$ \Term{constructible}{可构造}，进一步缩小 $U$ 后，$S\Phi|_U$ \Term{locally-constant}{局部常值}且\Term{constructible}{可构造}。
\end{proof}

\begin{SGAresult}{引理}{3.1.2}
设 $S$ 是\Term{locally-noetherian}{局部 Noether} \Term{scheme}{概形}，$f:X\to S$ 为\Term{morphism}{态射}，$F\to G$ 是 $X$ 上\Term{group-sheaf}{群层}的\Term{monomorphism}{单态射}。设 $\Psi$（分别地 $\Psi_1$）为 $F$-\Term{torsor}{挠子}（分别地 $G$-\Term{torsor}{挠子}）所成的叠，$\Phi=f_*\Psi$、$\Phi_1=f_*\Psi_1$。设给定 $S$ 上一个由\Term{finite-type}{有限型}\Term{etale}{平展} $S$-\Term{scheme}{概形} $T$ 表示的层，以及\Term{surjectivity}{满射}\Term{morphism}{态射}
\[a:T\longrightarrow S\Phi_1\simeq R^1f_*G,\]
使得 $X_T=X\times_S T$ 上存在群为 $G|_{X_T}$ 的\Term{torsor}{挠子} $Q$，它在 $R^1f_*G(T)$ 中定义 $T(T)$ 的恒等\Term{section}{截面}经 $a$ 所得的像。以 $f_T:X_T\to T$ 记典范\Term{morphism}{态射}，置 $F_T=F|_{X_T}$。若 $\Phi_1$ 是 $1$-\Term{constructible}{可构造}的，且 $f_{T*}(Q/F_T)$ \Term{constructible}{可构造}，则 $\Phi$ 为 $1$-\Term{constructible}{可构造}。
\end{SGAresult}
\SourcePage{318}
\begin{proof}
只需照搬 3.1.1 的证明，将存在\Term{morphism}{态射} $\Psi_i$ 这一点替换为存在\Term{isomorphism}{同构}
\[F'_i\xrightarrow{\sim}f_*(Q/F_T)|_{S'}.\]
\end{proof}

\begin{SGAresult}{注记}{3.1.3}
设 $k=k(s)$ 的特征为零。则 $\bar k$ 上维数 $\leq\dim X$ 的\Term{finite-type}{有限型}\Term{scheme}{概形}均可强消解奇点，且 3.1 的证明可给出下列结果：

a) 存在 $S$ 的一个非空开子集 $S_1$，使得对 $S_1$ 上每个极大点的特征为零的\Term{scheme}{概形} $S'$，以及 $X'=X\times_S S'$ 上每个\Term{locally-constant}{局部常值}\Term{constructible}{可构造}\Term{set-sheaf}{集合层} $F$，$(F,f')$ 相对于 $S'$ 在维数 $\leq0$ 内\Term{cohomologically-proper}{上同调固有}。

b) 若 $S$ 的全部\Term{residue-characteristic}{剩余特征}都是零，则存在 $S$ 的一个非空开子集 $S_1$，使得对 $S_1$ 上每个\Term{scheme}{概形} $S'$ 以及 $X'$ 上每个\Term{locally-constant}{局部常值}\Term{constructible}{可构造}\Term{group-sheaf}{群层} $F$，$(F,f')$ 相对于 $S'$ 在维数 $\leq1$ 内\Term{cohomologically-proper}{上同调固有}。
\end{SGAresult}
只需照搬 3.1 2) a) 的证明。3) 4 中所用的命题 2.7 也适用于\Term{locally-constant-sheaf}{局部常值层} $F$：沿用 3) 4 的记号，由于 $S$ 的所有\Term{residue-characteristic}{剩余特征}均为零，每个 $F$-\Term{torsor}{挠子}都在 $P$ 上相对于 $S$ \Term{tame-ramification}{温和分歧}。

\begin{SGAresult}{推论}{3.2}
设 $k$ 是特征 $p\geq0$ 的域，$p'$ 为不同于 $p$ 的素数组成的集合，$f:X\to k$ 为\Term{coherent-morphism}{凝聚态射}。

1) 对每个\Term{set-sheaf}{集合层} $F$，$(F,f)$ 在维数 $\leq0$ 内\Term{cohomologically-proper}{上同调固有}。

2) 设下列两个条件之一成立：

a) $f$ 为\Term{finite-type}{有限型}，且 $k$ 的一个\Term{algebraic-closure}{代数闭包}上维数 $\leq\dim X$ 的\Term{finite-type}{有限型}\Term{scheme}{概形}都可强消解奇点。

b) $k$ 的一个\Term{algebraic-closure}{代数闭包}上的\Term{finite-type}{有限型}\Term{scheme}{概形}都可强消解奇点。

则对每个 ind-$p'$-\Term{group-sheaf}{群层} $F$，$(F,f)$ 在维数 $\leq1$ 内\Term{cohomologically-proper}{上同调固有}。
\end{SGAresult}
\begin{proof}
设 $F$ 为\Term{set-sheaf}{集合层}（分别地 ind-$p'$-\Term{group-sheaf}{群层}）。由 SGA 4 IX 2.7.2，可将 $F$ 写成滤过\Term{inductive-limit}{归纳极限}
\[F=\varinjlim F_i,\]
其中 $F_i$ 是\Term{constructible}{可构造}\Term{set-sheaf}{集合层}（分别地\Term{constructible}{可构造} ind-$p'$-\Term{group-sheaf}{群层}）。由于 $f$ \Term{coherent}{凝聚}，$f_*$（分别地 $R^1f_*$）与\Term{inductive-limit}{归纳极限}可交换（SGA 4 VII 3.3）。若已知 $(F_i,f)$ 在维数 $\leq0$ 内\Term{cohomologically-proper}{上同调固有}（分别地每个\Term{set-sheaf}{集合层}都在维数 $\leq0$ 内\Term{cohomologically-proper}{上同调固有}，且 $(F_i,f)$ 在维数 $\leq1$ 内\Term{cohomologically-proper}{上同调固有}），则 $(F,f)$ 也如此。
\SourcePage{319}
故可设 $F$ \Term{constructible}{可构造}。若设 $f$ 为\Term{finite-type}{有限型}（分别地满足 a)），则结论由 3.1 1) b)（分别地 3.1 2) b)）得到。现不再假设 $f$ 为\Term{finite-type}{有限型}，来证明 3.2。对 $k$ 上每个\Term{scheme}{概形} $S'$ 和 $S'$ 的每个\Term{geometric-point}{几何点} $\bar s$，以 $\bar k$（分别地 $\overline{S'}$）记 $k$ 在 $\bar s$ 处（分别地 $S'$ 在 $\bar s$ 处）的\Term{strict-localization}{严格局部化}，$\bar X$ 为 $X$ 在 $\bar k$ 上的\Term{inverse-image}{逆像}，并考虑\Term{cartesian-square}{笛卡尔方块}
\[
\begin{tikzcd}
\bar X\arrow[d,"\bar f"']&\overline{X'}\arrow[l,"g"']\arrow[d,"\overline{f'}"]\\
\bar k&\overline{S'}\arrow[l].
\end{tikzcd}
\]
只需证明：对所有 $S'$ 和 $\bar s$，典范\Term{morphism}{态射}
\[
H^0(\bar X,\bar F)\longrightarrow H^0(\overline{X'},\overline{F'})
\quad\bigl(\text{分别地 }H^1(\bar X,\bar F)\longrightarrow H^1(\overline{X'},\overline{F'})\bigr)
\]
为\Term{isomorphism}{同构}。只需证明下列关系：
\[
\bar F\simeq g_*g^*\bar F
\qquad\bigl(\text{分别地 }R^1g_*(g^*\bar F)=0\bigr).\tag{*}
\]
在这一形式下，问题对 $X$ 的\Term{etale-topology}{平展拓扑}具有局部性。故可设 $X$ \Term{affine}{仿射}；利用取极限，又可设 $X$ 在 $k$ 上为\Term{finite-type}{有限型}。此时已知 $(F,f)$ 在维数 $\leq0$（分别地 $\leq1$）内\Term{cohomologically-proper}{上同调固有}，并且以 $X$ 上\Term{finite-type}{有限型}\Term{etale}{平展}\Term{scheme}{概形}替换 $X$ 后仍然如此，因而 (*) 得证。
\end{proof}

\begin{SGAresult}{定理}{3.3}
设 $S$ 是\Term{generic-point}{泛点}为 $s$ 的\Term{irreducible}{不可约}\Term{scheme}{概形}，$f:X\to S$ 为\Term{finite-presentation}{有限表示}\Term{morphism}{态射}。设 $k$ 的一个\Term{algebraic-closure}{代数闭包} $\bar k$ 上所有维数 $\leq\dim X_s$ 的\Term{finite-type}{有限型}\Term{scheme}{概形}都可消解奇点（EGA IV 7.9.1）。若 $L$ 为不同于 $S$ 的\Term{residue-characteristic}{剩余特征}的素数组成的集合，则存在 $S$ 的一个非空开子集 $S_1$，使 $f|_{S_1}$ 对 $L$ 是普遍\Term{locally-1-aspherical}{局部 $1$-无球}的。
\end{SGAresult}
\begin{proof}
可设 $S$ 整，$X$ \Term{reduced}{约化}（SGA 4 VIII 1.1）。利用取极限，可设 $S$ 为 Noether \Term{scheme}{概形}。此外，只需在\Term{finite}{有限}\Term{group-extension}{扩张} $S_1\to S$ 后证明定理，其中 $S_1$ 是\Term{integral-scheme}{整概形}，$S_1\to S$ 为\Term{etale}{平展}\Term{group-extension}{扩张}与\Term{radicial}{根性}\Term{surjectivity}{满射}\Term{group-extension}{扩张}的复合。

先证明缩小 $S$ 为非空开子集后，$f$ 普遍\Term{locally-0-acyclic}{局部 $0$-无环}。对 $k(s)$ 作\Term{radicial}{根性}\Term{group-extension}{扩张}后，可设\Term{morphism}{态射} $(X_s)_{\mathrm{red}}\to s$ \Term{separable}{可分}；于是缩小 $S$ 为非空开子集并作\Term{radicial}{根性}\Term{surjectivity}{满射}\Term{group-extension}{扩张}后，可设 $f$ \Term{flat}{平坦}且\Term{geometric-fibre}{几何纤维}\Term{separable}{可分}（EGA IV 12.1.1），这蕴含 $f$ 普遍\Term{0-acyclic}{$0$-无环}（SGA 4 XV 4.1）。

\SourcePage{320}
再证明缩小 $S$ 为非空开子集后，$f$ 对 $L$ 普遍\Term{locally-1-aspherical}{局部 $1$-无球}。对 $k(s)$ 作\Term{finite}{有限}\Term{group-extension}{扩张}后——这是允许的，因为它可视为\Term{etale}{平展}\Term{group-extension}{扩张}与\Term{radicial}{根性}\Term{group-extension}{扩张}的复合——可取\Term{proper}{固有}\Term{surjectivity}{满射} $p_s:Y_s\to X_s$，其中 $Y_s$ 是 $S$ 上与 $X_s$ 同维的\Term{smooth}{光滑}\Term{scheme}{概形}。还可设 $p_s$ 来自\Term{proper}{固有}\Term{surjectivity}{满射} $p:Y\to X$，其中 $Y$ 是 $S$ 上的\Term{smooth}{光滑}\Term{scheme}{概形}（EGA IV 9.6.1、12.1.6）。只需证明缩小 $S$ 为非空开子集后，对每个各方块均为\Term{cartesian-square}{笛卡尔方块}的图表
\[
\begin{tikzcd}
S''\arrow[d,"i"']&X''\arrow[l,"f''"']\arrow[d,"j"]\\
S'\arrow[d]&X'\arrow[l,"f'"']\arrow[d]\\
S&X\arrow[l,"f"],
\end{tikzcd}
\]
其中 $i$ \Term{etale}{平展}且\Term{finite-presentation}{有限表示}，以及 $S''$ 上每个 ind-$L$-\Term{group-sheaf}{群层} $F$，若 $\Phi$ 是 $F$-\Term{torsor}{挠子}所成的叠，则典范\Term{morphism}{态射}
\[f'^*i_*\Phi\longrightarrow j_*f''^*\Phi\]
是等价。置 $Z=Y\times_X Y$、$T=Y\times_X Y\times_X Y$，自然得到\Term{commutative-diagram}{交换图}
\begin{center}
\resizebox{\textwidth}{!}{\begin{tikzcd}[ampersand replacement=\&,column sep=small,row sep=large]
S''\arrow[d,"i"']\&X''\arrow[l,"f''"']\arrow[d,"j"]\&Y''\arrow[l,"p''"']\arrow[d]\&Z''\arrow[l,shift left=.45ex]\arrow[l,shift right=.45ex]\arrow[d]\&T''\arrow[l,shift left=.8ex]\arrow[l]\arrow[l,shift right=.8ex]\arrow[d]\\
S'\arrow[d]\&X'\arrow[l,"f'"']\arrow[d]\&Y'\arrow[l,"p'"']\arrow[d]\&Z'\arrow[l,shift left=.45ex]\arrow[l,shift right=.45ex]\arrow[d]\&T'\arrow[l,shift left=.8ex]\arrow[l]\arrow[l,shift right=.8ex]\arrow[d]\\
S\&X\arrow[l,"f"']\&Y\arrow[l,"p"']\&Z'\arrow[l,shift left=.45ex]\arrow[l,shift right=.45ex]\&T'\arrow[l,shift left=.8ex]\arrow[l]\arrow[l,shift right=.8ex].
\end{tikzcd}}
\end{center}
设 $q:Z\to X$、$r:T\to X$ 为典范\Term{morphism}{态射}；$q',r',q'',r''$ 的含义显然。由 1.11 2)，有下列本质交换、且各行\Term{exact}{正合}的图表：
\begin{center}
\resizebox{\textwidth}{!}{\begin{tikzcd}[ampersand replacement=\&,column sep=small,row sep=large]
f'^*i_*\Phi\arrow[r]\arrow[d,"a"']\&p'_*p'^*(f'^*i_*\Phi)\arrow[r,shift left=.45ex]\arrow[r,shift right=.45ex]\arrow[d,"b"]\&q'_*q'^*(f'^*i_*\Phi)\arrow[r,shift left=.8ex]\arrow[r]\arrow[r,shift right=.8ex]\arrow[d,"c"]\&r'_*r'^*(f'^*i_*\Phi)\arrow[d,"d"]\\
j_*f''^*\Phi\arrow[r]\&j_*p''_*p''^*(f''^*\Phi)\arrow[r,shift left=.45ex]\arrow[r,shift right=.45ex]\&j_*q''_*q''^*(f''^*\Phi)\arrow[r,shift left=.8ex]\arrow[r]\arrow[r,shift right=.8ex]\&j_*r''_*r''^*(f''^*\Phi).
\end{tikzcd}}
\end{center}
由于 $Y$ 在 $S$ 上\Term{smooth}{光滑}，\Term{morphism}{态射} $fp$ 对 $L$ 普遍\Term{locally-1-aspherical}{局部 $1$-无球}（SGA XV 2.1）；由 [2, VII 2.1.7]，$b$ 为\Term{category-equivalence}{范畴等价}。另一方面，缩小 $S$ 为非空开子集后，可设 $Z\to S$ 和 $T\to S$ 普遍\Term{locally-0-acyclic}{局部 $0$-无环}。
\SourcePage{321}
于是\Term{functor}{函子} $c,d$ \Term{fully-faithful}{完全忠实}；上图便说明 $a$ 为等价，证明完成。
\end{proof}

\begin{SGAresult}{推论}{3.4}
设 $k$ 是特征 $p\geq0$ 的域，$p'$ 为不同于 $p$ 的素数组成的集合，$f:X\to k$ 为\Term{coherent-morphism}{凝聚态射}。设下列两个条件之一成立：

a) $f$ 为\Term{finite-type}{有限型}，且 $k$ 的一个\Term{algebraic-closure}{代数闭包}上维数 $\leq\dim X$ 的\Term{finite-type}{有限型}\Term{scheme}{概形}都可消解奇点。

b) $k$ 的一个\Term{algebraic-closure}{代数闭包}上的\Term{finite-type}{有限型}\Term{scheme}{概形}都可消解奇点。

则 $f$ 对 $p'$ 普遍\Term{locally-1-aspherical}{局部 $1$-无球}。
\end{SGAresult}
情形 a) 由 3.3 得到。情形 b) 中，问题对 $X$ 具有局部性，可设 $X$ \Term{affine}{仿射}；利用取极限（SGA 4 XV 1.3），归约到 $X$ 在 $k$ 上\Term{finite-type}{有限型}的情形。

\begin{SGAresult}{推论}{3.5}
设 $S$ 是\Term{generic-point}{泛点}为 $s$ 的\Term{irreducible}{不可约}\Term{scheme}{概形}，$f:X\to S$ 为\Term{finite-presentation}{有限表示}\Term{morphism}{态射}。设 $k(s)$ 的一个\Term{algebraic-closure}{代数闭包} $\bar k$ 上所有维数 $\leq\dim X_s$ 的\Term{finite-type}{有限型}\Term{scheme}{概形}都可强消解奇点（SGA 5 I 3.1.5）。若 $L$ 为不同于 $S$ 的\Term{residue-characteristic}{剩余特征}的素数组成的集合，则存在 $S$ 的一个非空开子集 $S_1$，使对 $S_1$ 的\Term{geometric-point}{几何点}的每个\Term{specialization}{特殊化} $\bar s_1\to\bar s_2$，\Term{specialization-morphism}{特殊化态射}（2.10）
\[\pi_1^L(X_{\bar s_1})\longrightarrow\pi_1^L(X_{\bar s_2})\]
都是\Term{bijection}{双射}。
\end{SGAresult}
由 3.1、3.3，缩小 $S$ 为非空开子集后，可设 $f$ 对 $L$ \Term{locally-1-aspherical}{局部 $1$-无球}，并且对 $X$ 上每个\Term{finite}{有限}常值 $L$-\Term{group-sheaf}{群层} $F$，$(F,f)$ 在维数 $\leq1$ 内\Term{cohomologically-proper}{上同调固有}。由 1.14，对 $S_1$ 的\Term{geometric-point}{几何点}的每个\Term{specialization}{特殊化} $\bar s_1\to\bar s_2$，\Term{specialization-morphism}{特殊化态射}
\[(R^1f_*F)_{\bar s_2}\longrightarrow(R^1f_*F)_{\bar s_1}\]
均为\Term{bijection}{双射}。本推论只是将上述结论用\Term{fundamental-group}{基本群}的语言表述出来。

\SGAsection{4}{同伦正合列}
\textbf{4.0.} 设 $X,S$ 是两个\Term{connected}{连通}\Term{scheme}{概形}，$f:X\to S$ 为\Term{morphism}{态射}，$a$ 是 $X$ 的\Term{geometric-point}{几何点}，$L$ 为一个素数集合。设 $K$ 是典范\Term{homomorphism}{同态} $\pi_1(X,a)\to\pi_1(S,a)$ 的核，$N$ 是 $K$ 的最小\Term{normal}{正规}投子群，使 $K/N$ 为一个投-$L$-群 $K^L$。
\SourcePage{322}
则 $N$ 在 $\pi_1(X,a)$ 中\Term{normal}{正规}，以
\[\pi'_1(X,a)\]
记 $\pi_1(X,a)$ 对 $N$ 的商。若 $a$ 是\Term{geometric-fibre}{几何纤维} $X_{\bar s}$ 的一个\Term{geometric-point}{几何点}，则典范\Term{morphism}{态射}
\[\pi_1(X_{\bar s},a)\longrightarrow\pi_1(X,a)\longrightarrow\pi_1(S,a)\]
给出典范\Term{morphism}{态射}
\[\pi_1^L(X_{\bar s},a)\xrightarrow{u}\pi'_1(X,a)\xrightarrow{v}\pi_1(S,a).\]
有 $vu=0$。

\begin{SGAresult}{命题}{4.1}
设 $S$ 为\Term{connected}{连通}\Term{scheme}{概形}，$f:X\to S$ 为\Term{locally-0-acyclic}{局部 $0$-无环}\Term{morphism}{态射}（SGA 4 XV 1.11），并设 $f$ 也 \Term{0-acyclic}{$0$-无环}；当 $f$ \Term{coherent}{凝聚}时，这等价于 $f$ 的\Term{geometric-fibre}{几何纤维}\Term{connected}{连通}（SGA 4 XV 1.16）。设 $L$ 为一个素数集合。若 $S'$ 是 $S$ 上的\Term{etale}{平展}\Term{scheme}{概形}，以 $X',f'$ 记 $X,f$ 在 $S'$ 上的\Term{inverse-image}{逆像}。设对 $S$ 的每个\Term{etale-covering}{平展覆盖} $S'$，以及 $X'$ 的每个\Term{etale-covering}{平展覆盖} $E$，只要 $E$ 是一个以 $L$-群为群的 \Term{galois-covering}{Galois 覆盖}的商，就有 $(E,f')$ 相对于 $S'$ 在维数 $\leq0$ 内\Term{cohomologically-proper}{上同调固有}，且 $f'_*E$ \Term{constructible}{可构造}。则若 $\bar s$ 是 $S$ 的\Term{geometric-point}{几何点}，$a$ 是\Term{fiber}{纤维} $X_{\bar s}$ 的\Term{geometric-point}{几何点}，\Term{group-homomorphism}{群同态}列
\[
\pi_1^L(X_{\bar s},a)\xrightarrow{u}\pi'_1(X,a)\xrightarrow{v}\pi_1(S,a)\longrightarrow1\tag{4.1.1}
\]
\Term{exact}{正合}。
\end{SGAresult}
\begin{proof}
本命题推广了 X 1.4，其证明也照搬该处。

先证 $v$ \Term{surjectivity}{满射}。只需证明，对 $S$ 的每个\Term{connected}{连通}\Term{etale-covering}{平展覆盖} $S'$，$X'$ 也\Term{connected}{连通}（V 6.9）。设集合 $C$ 至少含两个元素。由于 $f$ 为 \Term{0-acyclic}{$0$-无环}，典范\Term{morphism}{态射}
\[H^0(S',C_{S'})\longrightarrow H^0(X',C_{X'})\]
是\Term{bijection}{双射}，故 $X'$ \Term{connected}{连通}，从而 $v$ \Term{surjectivity}{满射}。

由 $K^L$ 的定义（4.0），有\Term{exact-sequence}{正合列}
\[1\longrightarrow K^L\longrightarrow\pi'_1(X,a)\longrightarrow\pi_1(S,a)\longrightarrow1.\]
设 $\widetilde S$ 为 $S$ 的\Term{universal-covering}{万有覆盖}，$\widetilde X=\widetilde S\times_S X$。群 $K^L$ 分类这样一些以 $L$-群为群的 \Term{galois-covering}{Galois 覆盖} $P$：存在 $S$ 的\Term{etale-covering}{平展覆盖} $S'$ 及 $X'=X\times_S S'$ 的 \Term{galois-covering}{Galois 覆盖} $Q$，使有\Term{isomorphism}{同构} $P\simeq Q\times_{X'}\widetilde X$。列 (4.1.1) \Term{exact}{正合}，当且仅当典范\Term{morphism}{态射}
\[\pi_1^L(X_{\bar s},a)\longrightarrow K^L\]
\Term{surjectivity}{满射}。由 $K^L$ 的解释，这又等价于：对 $S$ 的每个\Term{etale-covering}{平展覆盖} $S'$ 和 $X'$ 的每个以 $L$-群为群的 \Term{galois-covering}{Galois 覆盖} $Q$，若 $P=Q\times_{X'}\widetilde X$ \Term{connected}{连通}，则 $Q|_{X_{\bar s}}$ \Term{connected}{连通}。
\SourcePage{323}
来证明这一条件成立。设 $S'$ 是 $S$ 的\Term{etale-covering}{平展覆盖}，$Q$ 为 $X'$ 的群为 $L$-群 $F$ 的 \Term{galois-covering}{Galois 覆盖}，并设 $Q|_{X_{\bar s}}$ 不\Term{connected}{连通}；要证明以 $S'$ 的一个\Term{etale-covering}{平展覆盖}替换 $S'$ 后，$Q$ 不\Term{connected}{连通}。存在 $F$ 的一个真子群 $G$ 及 $G|_{X_{\bar s}}$ 的一个\Term{torsor}{挠子} $R$，使 $Q|_{X_{\bar s}}$ 由 $R$ 沿 $G\to F$ 作结构群\Term{group-extension}{扩张}得到。$X'$ 的\Term{etale-covering}{平展覆盖} $E=Q/G$ 在 $X_{\bar s}$ 上有\Term{section}{截面}。由 1.16，$f_*E$ \Term{locally-constant}{局部常值}且\Term{constructible}{可构造}；以 $S'$ 的一个\Term{etale-covering}{平展覆盖}替换 $S'$ 后，甚至可设 $f_*E$ 常值。又由于 $(E,f')$ 相对于 $S'$ 在维数 $\leq0$ 内\Term{cohomologically-proper}{上同调固有}，且 $H^0(X_{\bar s},E|_{X_{\bar s}})$ 非空，$E$ 有\Term{section}{截面}。这说明 $Q$ 不\Term{connected}{连通}，证明完成。
\end{proof}
由 4.1 得到下列引理，4.6 将用到它。

\begin{SGAresult}{引理}{4.2}
设 $S$ 为\Term{connected}{连通}\Term{scheme}{概形}，$f:X\to S$ 为\Term{locally-0-acyclic}{局部 $0$-无环}且 \Term{0-acyclic}{$0$-无环}的\Term{morphism}{态射}，$L$ 为一个素数集合。设对 $X$ 上每个\Term{finite}{有限}常值 $L$-\Term{group-sheaf}{群层} $F$，$(F,f)$ 相对于 $S$ 在维数 $\leq1$ 内\Term{cohomologically-proper}{上同调固有}；并设对 $S$ 的每个\Term{etale-covering}{平展覆盖} $S'$ 和 $X'$ 的每个\Term{etale-covering}{平展覆盖} $E$，只要 $E$ 是一个以 $L$-群为群的 \Term{galois-covering}{Galois 覆盖}的商，就有 $f'_*E$ \Term{constructible}{可构造}。则若 $\bar s$ 是 $S$ 的\Term{geometric-point}{几何点}，$a$ 是\Term{fiber}{纤维} $X_{\bar s}$ 的\Term{geometric-point}{几何点}，\Term{group-homomorphism}{群同态}列
\[\pi_1^L(X_{\bar s},a)\longrightarrow\pi'_1(X,a)\longrightarrow\pi_1(S,a)\longrightarrow1\]
\Term{exact}{正合}。
\end{SGAresult}
4.1 的假设成立。事实上，由 1.13 3)，对 $S$ 上每个\Term{etale}{平展}\Term{scheme}{概形} $S'$ 以及 $X'$ 的每个\Term{etale-covering}{平展覆盖} $E$，只要 $E$ 是一个以 $L$-群为群的 \Term{galois-covering}{Galois 覆盖}的商，就有 $(E,f')$ 相对于 $S'$ 在维数 $\leq0$ 内\Term{cohomologically-proper}{上同调固有}。

\begin{SGAresult}{命题}{4.3}
设 $S$ 为\Term{connected}{连通}\Term{scheme}{概形}，$L$ 为一个素数集合，$f:X\to S$ 为 \Term{0-acyclic}{$0$-无环}且对 $L$ \Term{locally-1-aspherical}{局部 $1$-无球}的\Term{morphism}{态射}（SGA 4 XV 1.11），$g:S\to X$ 为 $f$ 的一个\Term{section}{截面}。设 $\bar s$ 是 $S$ 的\Term{geometric-point}{几何点}，$a$ 是\Term{fiber}{纤维} $X_{\bar s}$ 的\Term{geometric-point}{几何点}。设对每个常值 $L$-\Term{group-sheaf}{群层} $F$，$(F,f)$ 在维数 $\leq1$ 内\Term{cohomologically-proper}{上同调固有}，$F$-\Term{torsor}{挠子}所成的叠经 $f$ 的\Term{direct-image}{正像}是 $1$-\Term{constructible}{可构造}叠（0）；并设对 $X'$ 的每个\Term{etale-covering}{平展覆盖} $E$，只要 $E$ 是一个以 $L$-群为群的 \Term{galois-covering}{Galois 覆盖}的商，就有 $f'_*E$ \Term{constructible}{可构造}。则\Term{group-homomorphism}{群同态}列
\[
1\longrightarrow\pi_1^L(X_{\bar s},a)\xrightarrow{u}\pi'_1(X,a)\xrightarrow{v}\pi_1(S,a)\longrightarrow1\tag{4.3.1}
\]
\Term{exact}{正合}。
\end{SGAresult}
\SourcePage{324}
\begin{proof}
结合 4.2，只需证明 $u$ 的\Term{injectivity}{单射}性；亦即证明：对 $X_{\bar s}$ 的每个以 $L$-群 $C$ 为群的\Term{principal-covering}{主覆盖} $\bar Z$，存在 $X$ 的\Term{etale-covering}{平展覆盖} $Z$，以及从 $Z|_{X_{\bar s}}$ 的某个\Term{connected-component}{连通分支}到 $\bar Z$ 的\Term{morphism}{态射}（V 6.8）。设 $\bar Z$ 为这样的\Term{principal-covering}{主覆盖}，$\bar z$ 为其在 $H^1(X_{\bar s},C_{X_{\bar s}})$ 中的类。由 1.5 d)，有典范\Term{isomorphism}{同构}
\[(R^1f_*C_X)_{\bar s}\xrightarrow{\sim}H^1(X_{\bar s},C_{X_{\bar s}}),\]
由 1.16，$R^1f_*C_X$ 是\Term{locally-constant}{局部常值}\Term{constructible}{可构造}层。因此存在 $S$ 的\Term{etale-covering}{平展覆盖} $S'$，使 $R^1f_*C_X|_{S'}$ 常值。若 $\bar s\to S'$ 是位于\Term{geometric-point}{几何点} $\bar s\to S$ 上方的\Term{geometric-point}{几何点}，则 $H^0(S',R^1f_*C_X)$ 中存在元素 $z$，其在 $H^1(X_{\bar s},C_{X_{\bar s}})$ 中的像为 $\bar z$。由下文引理 4.3.1，存在 $S'$ 的\Term{etale-covering}{平展覆盖} $S'_1$，以及 $X'_1$ 上群为 $C$ 的\Term{torsor}{挠子} $P$，使它在 $H^0(S'_1,R^1f_*C)$ 中的像等于 $z$ 的限制。\Term{torsor}{挠子} $P$ 由 $X'_1=X\times_S S'_1$ 的一个\Term{etale-covering}{平展覆盖} $Z$ 表示，且 $Z\times_{X'_1}X_{\bar s}$ \Term{isomorphism}{同构}于 $\bar Z$。把 $Z$ 视为 $X$ 的\Term{etale-covering}{平展覆盖}，即得 $Z\times_X X_{\bar s}\to\bar Z$ 的\Term{morphism}{态射}，证明完成。
\end{proof}

\begin{SGAresult}{引理}{4.3.1}
设 $f:X\to S$ 为 \Term{0-acyclic}{$0$-无环}且\Term{locally-0-acyclic}{局部 $0$-无环}的\Term{morphism}{态射}，$g$ 为 $f$ 的一个\Term{section}{截面}。设 $C$ 是\Term{finite}{有限}常值群，满足 $(C_X,f)$ 在维数 $\leq0$ 内\Term{cohomologically-proper}{上同调固有}，且 $C_X$-\Term{torsor}{挠子}所成的叠经 $f$ 的\Term{direct-image}{正像}\Term{constructible}{可构造}。则对 $H^0(S,R^1f_*C_X)$ 的每个\Term{section}{截面} $z$，存在 $S$ 的一个\Term{etale-covering}{平展覆盖} $S_1$；置 $X_1=X\times_S S_1$ 后，还存在 $H^1(X_1,C_{X_1})$ 中的一个元素，其沿典范\Term{morphism}{态射}
\[H^1(X_1,C_{X_1})\longrightarrow H^0(S_1,R^1f_*C_{X_1})\]
的像等于 $z$ 在 $H^0(S_1,R^1f_*C_{X_1})$ 上的限制。
\end{SGAresult}
\begin{proof}
对 $S$ 上每个\Term{etale}{平展}\Term{scheme}{概形} $S'$，置 $X'=X\times_S S'$，并以 $g'$（分别地 $F'$ 等）记 $g$（分别地 $F$ 等）沿 $S'\to S$ 的\Term{inverse-image}{逆像}。$S$ 上由
\[
\begin{aligned}
G(S')=\{&X'\text{ 上群为 }C_{X'}\text{、配备同构 }g'^*P\xrightarrow{\sim}C_{S'}\\
&\text{的挠子 }P\text{ 的同构类}\}
\end{aligned}
\]
定义的\Term{presheaf}{预层} $G$ 是层。事实上，由\Term{descent}{下降}，这归结于：$X'$ 上\Term{torsor}{挠子} $P$ 的一个\Term{isomorphism}{同构}由其在 $g'(S')$ 上的限制唯一确定。此外，有\Term{surjectivity}{满射}\Term{morphism}{态射}
\[G\longrightarrow R^1f_*F.\]
设 $z\in H^0(S,R^1f_*C_X)$，$H$ 为 $G$ 中 $z$ 的\Term{inverse-image}{逆像}\Term{subsheaf}{子层}。只需证明 $H$ \Term{locally-constant}{局部常值}且\Term{constructible}{可构造}。此性质对 $S$ 具有局部性，故可设 $z$ 来自 $H^1(X,C_X)$ 的一个元素，由\Term{torsor}{挠子} $P$ 表示，且 $g^*P\simeq C_S$。给定\Term{isomorphism}{同构} $i:g^*P\xrightarrow{\sim}C_S$，等价于给定 $\underline{\operatorname{Aut}}_{C_S}(g^*P)$ 的一个整体\Term{section}{截面}；两个\Term{isomorphism}{同构} $i,i'$ 在 $G(X)$ 中定义同一个元素，当且仅当 $ii'^{-1}$ 是 $\operatorname{Aut}_{C_X}(P)$ 中一个元素的像。
\SourcePage{325}
考虑典范\Term{injectivity}{单射}
\[f_*\underline{\operatorname{Aut}}_{C_X}(P)\longrightarrow\underline{\operatorname{Aut}}_{C_S}(g^*P)\simeq C,\]
便可将 $H$ 识别为 $\underline{\operatorname{Aut}}_{C_S}(g^*P)$ 对 $f_*\underline{\operatorname{Aut}}_{C_X}(P)$ 的商。由 1.16，$f_*\underline{\operatorname{Aut}}_{C_X}(P)$ \Term{locally-constant}{局部常值}，故 $H$ 亦如此，证明完成。
\end{proof}

\begin{SGAresult}{例}{4.4}
注意：若 $S$ 是\Term{connected}{连通}\Term{scheme}{概形}，当 $f$ 为\Term{proper}{固有}、\Term{flat}{平坦}、\Term{finite-presentation}{有限表示}\Term{morphism}{态射}且\Term{geometric-fibre}{几何纤维}\Term{separable}{可分}并\Term{connected}{连通}时，4.1 的假设成立，而 $L$ 可任意（参见 X 1.3）。若此外 $f$ \Term{smooth}{光滑}且有\Term{section}{截面}，并以 $L$ 记不同于 $S$ 的\Term{residue-characteristic}{剩余特征}的素数组成的集合，则 4.3 的假设成立（SGA 4 XV 2.1、XVI 5.2）。

若 $S$ \Term{connected}{连通}，存在\Term{proper}{固有}\Term{finite-presentation}{有限表示}的 $S$-\Term{scheme}{概形} $Z$，它在 $S$ 上\Term{flat}{平坦}且\Term{geometric-fibre}{几何纤维}\Term{separable}{可分}并\Term{connected}{连通}，而 $X$ 是 $Z$ 中一个相对于 $S$ 具有\Term{normal-crossings}{正规交叉}的\Term{divisor}{除子}的补集，并取 $L$ 为不同于 $S$ 的\Term{residue-characteristic}{剩余特征}的素数组成的集合，则 4.1 的假设也成立（2.9）。若此外 $f$ \Term{smooth}{光滑}且有\Term{section}{截面}，4.3 的假设也成立。
\end{SGAresult}

\textbf{4.5.} 沿用 4.3 的记号和假设。若 $\bar s$ 是 $S$ 的\Term{geometric-point}{几何点}，$a=g(\bar s)$，则\Term{section}{截面} $g$ 给出\Term{morphism}{态射}
\[w:\pi_1(S,a)\longrightarrow\pi'_1(X,a),\]
从而 $\pi'_1(X,a)$ 识别为 $\pi_1(X_{\bar s},a)$ 与 $\pi_1(S,a)$ 的\Term{semidirect-product}{半直积}。因此\Term{profinite}{投有限}群 $\pi_1(S,a)$ 作用于 $\pi_1(X_{\bar s},a)$。由于 $\pi_1^L(X_{\bar s},a)$ 是在 $\pi_1(S,a)$ 的作用下不变的群的严格\Term{projective-limit}{投射极限}，给定带此作用的 $\pi_1^L(X_{\bar s},a)$，等价于给定 $S$ 上\Term{finite}{有限}\Term{etale}{平展}\Term{group-scheme}{群概形}的一个严格\Term{inverse-system}{投射系}，记为
\[\pi_1^L(X/S,g,\bar s)\qquad\text{或简记为 }\pi_1^L(X/S,g).\]
有如下性质。

\textbf{4.5.1.} 对 $S$ 上每个\Term{fiber}{纤维}为 $L$-群的\Term{finite}{有限}\Term{etale}{平展}\Term{group-scheme}{群概形} $G$，设 $E$ 为 $G$ 在 $X$ 上的\Term{inverse-image}{逆像} $G_X$ 的、配备\Term{isomorphism}{同构} $g^*P\xrightarrow{\sim}G$ 的\Term{torsor}{挠子} $P$ 的类所成的集合，则 $E$ 典范\Term{isomorphism}{同构}于集合
\[\operatorname{Hom}_S(\pi_1^L(X/S,g,\bar s),G)\big/\text{$G$ 的内自同构}.\]

\SourcePage{326}
\textbf{4.5.2.} 对 $S$ 上每个\Term{fiber}{纤维}为 $L$-群的\Term{finite}{有限}\Term{etale}{平展}\Term{group-scheme}{群概形} $G$，层 $R^1f_*G_X$ 典范\Term{isomorphism}{同构}于下列\Term{presheaf}{预层}的相伴层：
\[S'\longmapsto\operatorname{Hom}_{S'}(\pi_1^L(X/S,g,\bar s),G)\big/\text{$G$ 的内自同构},\]
其中 $S'$ 为 $S$ 上的\Term{etale}{平展}\Term{scheme}{概形}。

\textbf{4.5.3.} 设 $S'$ 是\Term{connected}{连通} $S$-\Term{scheme}{概形}，$\bar s$ 为 $S'$ 的\Term{geometric-point}{几何点}，$X',g'$ 分别为 $X,g$ 在 $S'$ 上的\Term{inverse-image}{逆像}。则 $\pi_1^L(X'/S',g',\bar s)$ 典范\Term{isomorphism}{同构}于 $\pi_1^L(X/S,g,\bar s)$ 在 $S'$ 上的\Term{inverse-image}{逆像}。对 $S$ 的每个\Term{geometric-point}{几何点} $\xi$，\Term{fiber}{纤维} $\pi_1^L(X/S,g,\bar s)_\xi$ \Term{isomorphism}{同构}于 $\pi_1^L(X_\xi)$。

事实上，给定 $G$ 等价于给定一个带 $\pi_1(S,a)$ 作用的抽象 $L$-群 $\mathbf G$，从而得到 $\pi'_1(X,a)$ 在 $\mathbf G$ 上的作用。4.5.1 的\Term{isomorphism}{同构}来自典范\Term{morphism}{态射}
\[
H^1(\pi'_1(X,a),\mathbf G)\longrightarrow
H^1(\pi_1^L(X_{\bar s},a),\mathbf G)
=\operatorname{Hom}(\pi_1^L(X_{\bar s},a),\mathbf G)/\operatorname{Int}(G)
\]
在子集 $E$ 上的限制；$E$ \Term{bijection}{双射}到 $\pi_1^L(X_{\bar s},a)\to\mathbf G$ 中与 $\pi_1(S,a)$ 的作用相容的\Term{morphism}{态射}所成的子集。断言 4.5.3 由 $\pi_1^L(X/S,g,\bar s)$ 的定义结合\Term{homotopy-exact-sequence}{同伦正合列} (4.3.1) 得到；4.5.2 则由 4.5.1、4.5.3 推出。

\begin{SGAresult}{命题（Künneth 公式）}{4.6}
设 $k$ 是特征 $p\geq0$ 的\Term{separably-closed}{可分闭}域，$X,Y$ 是两个\Term{connected}{连通} $k$-\Term{scheme}{概形}，$a$ 为 $X$ 的\Term{geometric-point}{几何点}，$b$ 为 $Y$ 的\Term{geometric-point}{几何点}，$c$ 为 $X\times_kY$ 中位于 $a,b$ 上方的\Term{geometric-point}{几何点}。设下列两个条件之一成立：

a) $X$ 在 $k$ 上为\Term{finite-type}{有限型}，且 $k$ 的一个\Term{algebraic-closure}{代数闭包} $\bar k$ 上维数 $\leq\dim X$ 的\Term{finite-type}{有限型}\Term{scheme}{概形}都可强消解奇点（SGA 5 I 3.1.5）。

b) $X$ \Term{quasi-compact}{拟紧}且\Term{quasi-separated}{拟分离}，且 $\bar k$ 上每个\Term{finite-type}{有限型}\Term{scheme}{概形}都可强消解奇点。

若 $p'$ 为不同于 $p$ 的素数组成的集合，则从投影
\[X\times_kY\longrightarrow X,\qquad X\times_kY\longrightarrow Y\]
所诱导的\Term{fundamental-group}{基本群}\Term{homomorphism}{同态}得到的\Term{morphism}{态射}
\[\pi_1^{p'}(X\times_kY,c)\longrightarrow\pi_1^{p'}(X,a)\times\pi_1^{p'}(Y,b)\tag{4.6.0}\]
是\Term{isomorphism}{同构}。
\end{SGAresult}
\begin{proof}
可设 $k$ \Term{algebraically-closed}{代数闭}，$X$ \Term{reduced}{约化}（SGA 4 VIII 1.1）。设 $Z=X\times_kY$，$g:X\to k$、$f:Z\to Y$ 为典范\Term{morphism}{态射}。由 3.5，\Term{morphism}{态射} $g$，因而 $f$，对 $p'$ 普遍\Term{locally-1-aspherical}{局部 $1$-无球}。由于 $X$ \Term{connected}{连通}，$f$ 为 \Term{0-acyclic}{$0$-无环}（SGA 4 XV 1.16）。另一方面，由 3.2，对每个\Term{finite}{有限} $p'$-群 $C$，$(C_X,g)$ 相对于 $k$ 在维数 $\leq1$ 内\Term{cohomologically-proper}{上同调固有}。
\SourcePage{327}
故 $(C_Z,f)$ 相对于 $Y$ 在维数 $\leq1$ 内\Term{cohomologically-proper}{上同调固有}（1.5 c)），而 $f_*C_Z$、$R^1f_*C_Z$ 是\Term{constant-sheaf}{常值层}。因此 $g$ 满足 4.2 的全部假设，得到\Term{exact-sequence}{正合列}
\[\pi_1^{p'}(X_b,c)\longrightarrow\pi_1^{p'}(Z,c)\longrightarrow\pi_1^{p'}(Y,b)\longrightarrow1.\]
此外，复合\Term{morphism}{态射}
\[\pi_1^{p'}(X_b,c)\longrightarrow\pi_1^{p'}(Z,c)\longrightarrow\pi_1^{p'}(Y,b)\]
为\Term{isomorphism}{同构}，故有\Term{exact-sequence}{正合列}
\[1\longrightarrow\pi_1^{p'}(X,a)\longrightarrow\pi_1^{p'}(Z,c)\longrightarrow\pi_1^{p'}(Y,b)\longrightarrow1.\]
另一方面，\Term{morphism}{态射} (4.6.0) 给出从此\Term{exact-sequence}{正合列}到\Term{exact-sequence}{正合列}
\[1\longrightarrow\pi_1^{p'}(X,a)\longrightarrow
\pi_1^{p'}(X,a)\times\pi_1^{p'}(Y,b)\times\pi_1^{p'}(Y,b)\longrightarrow1\]
的\Term{morphism}{态射}，从而 (4.6.0) 为\Term{isomorphism}{同构}。
\end{proof}

\textbf{4.7.} 设
\[
\begin{tikzcd}
X\arrow[d,"f"']&X_U\arrow[l]\arrow[d,"f_U"]&X_{\bar s}\arrow[l]\arrow[d]\\
S&U\arrow[l]&\bar s\arrow[l]
\end{tikzcd}
\]
是各方块均为\Term{cartesian-square}{笛卡尔方块}的图表，其中 $S$ 是\Term{arcwise-connected}{道路连通}\Term{scheme}{概形}（SGA 4 IX 2.12），$U$ 为 $S$ 的\Term{connected}{连通}开子集，$\bar s$ 为 $U$ 的\Term{geometric-point}{几何点}。设 $a$ 为 $X_{\bar s}$ 的\Term{geometric-point}{几何点}，$L$ 为素数集合。设 $g$ 为 $f$ 的\Term{section}{截面}，并设下列条件成立：

a) \Term{morphism}{态射} $f$ 为 \Term{0-acyclic}{$0$-无环}且\Term{locally-0-acyclic}{局部 $0$-无环}；对 $S$ 的每个\Term{etale-covering}{平展覆盖} $S'$ 和 $X\times_S S'$ 的每个\Term{etale-covering}{平展覆盖} $E$，只要 $E$ 是一个以 $L$-群为群的 \Term{galois-covering}{Galois 覆盖}的商，就有 $(E,f_{(S')})$ 相对于 $S'$ 在维数 $\leq0$ 内\Term{cohomologically-proper}{上同调固有}。

b) \Term{morphism}{态射} $f_U$ 对 $L$ \Term{locally-1-aspherical}{局部 $1$-无球}；对 $X_U$ 上每个\Term{finite}{有限}常值 $L$-\Term{group-sheaf}{群层} $F$，$(F,f_U)$ 在维数 $\leq1$ 内\Term{cohomologically-proper}{上同调固有}，且 $R^1f_{U*}F$ 的\Term{fiber}{纤维}\Term{finite}{有限}。

由 4.1、4.3 得到下列\Term{commutative-diagram}{交换图}，其各行\Term{exact}{正合}：
\begin{center}
\resizebox{.88\textwidth}{!}{\begin{tikzcd}[ampersand replacement=\&,column sep=small]
1\arrow[r]\&\pi_1^L(X_{\bar s},a)\arrow[r]\arrow[d,equal]\&\pi'_1(X_U,a)\arrow[r]\arrow[d]\&\pi_1(U,a)\arrow[r]\arrow[d]\&1\\
\&\pi_1^L(X_{\bar s},a)\arrow[r]\&\pi'_1(X,a)\arrow[r]\&\pi_1(S,a)\arrow[r]\&1.
\end{tikzcd}}\quad(4.7.0)
\end{center}
利用\Term{section}{截面} $g$，有\Term{morphism}{态射}
\[\pi_1(U,a)\longrightarrow\pi'_1(X_U,a),\qquad\pi_1(S,a)\longrightarrow\pi'_1(X,a).\]
\SourcePage{328}
由此得到从 $\pi_1(S,a)$ 与 $\pi'_1(X_U,a)$ 在 $\pi_1(U,a)$ 上的\Term{amalgamated-sum}{合并和}到 $\pi'_1(X,a)$ 的\Term{morphism}{态射}：
\[\varphi:\pi=\pi_1(S,a)\underset{\pi_1(U,a)}{\amalg}\pi'_1(X_U,a)\longrightarrow\pi'_1(X,a).\tag{4.7.1}\]
再设下列条件成立：

c) 若 $T=S-U$，则 $\operatorname{prof\text{ét}}_T(S)\geq2$（SGA 2 XIV 1.1）。

则将 $S$ 的\Term{etale-covering}{平展覆盖}对应到其在 $U$ 上的限制的\Term{functor}{函子}\Term{fully-faithful}{完全忠实}（SGA 2 XVI 1.4）。因此，\Term{morphism}{态射} $\pi_1(U,a)\to\pi_1(S,a)$ \Term{surjectivity}{满射}（V 6.9）；由图 (4.7.0)，$\pi'_1(X_U,a)\to\pi'_1(X,a)$ 也\Term{surjectivity}{满射}，故 $\varphi$ 更是\Term{epimorphism}{满态射}。置
\[K=\operatorname{Ker}(\pi_1(U,a)\longrightarrow\pi_1(S,a)).\tag{4.7.2}\]
(4.7.1) 中的群 $\pi$ 识别为 $\pi'_1(X_U,a)$ 对 $K$ 在其中的像所生成的闭\Term{normal-subgroup}{正规子群}的商。以 $L$ 记这一像，将 $\pi'_1(X_U,a)$ 视为 $\pi_1^L(X_{\bar s},a)$ 与 $\pi_1(U,a)$ 的\Term{semidirect-product}{半直积}。于是 $K$ 通过\Term{inner-automorphism}{内自同构}作用于 $\pi_1^L(X_{\bar s},a)$，而商群 $\pi=\pi'_1(X_U,a)/L$ 识别为 $\pi_1^L(X_{\bar s},a)$ 在 $K$ 下的\Term{coinvariants}{余不变量}群 $\pi_1^L(X_{\bar s},a)_K$ 与 $\pi_1(U,a)/K=\pi_1(S,a)$ 的\Term{semidirect-product}{半直积}。最终得到\Term{epimorphism}{满态射}
\[\varphi:\pi=\pi_1^L(X_{\bar s},a)_K\cdot\pi_1(S,a)\longrightarrow\pi'_1(X,a).\tag{4.7.3}\]
下面的命题给出 $\varphi$ 为\Term{isomorphism}{同构}的条件。

\begin{SGAresult}{命题}{4.7.4}
沿用 4.7 的记号。除条件 a)、b)、c) 外，再设下列条件成立：

d) 对 $T=S-U$ 的每个点 $t$，\Term{morphism}{态射} $f$ 在 $g(t)$ 处对 $L$ \Term{locally-1-aspherical}{局部 $1$-无球}。

e) 对每个点 $t\in T$，\Term{fiber}{纤维} $X_t$ 的每个\Term{irreducible-component}{不可约分支}都包含 $g(t)$；并且对 $X_t-\{g(t)\}$ 的每个非极大点 $x$，有
\[\operatorname{prof\,hop}_x(X)\geq3\qquad\text{（SGA 2 XIV 1.2）},\]
且环 $\mathcal O_{X,x}$ 为 Noether 环。

则\Term{morphism}{态射} (4.7.3) 是\Term{isomorphism}{同构}。
\end{SGAresult}
\begin{proof}
如前所述，群 $\pi$ 识别为 $\pi'_1(X_U,a)$ 对 $K$（4.7.2）在其中的像所生成的闭\Term{normal-subgroup}{正规子群} $L$ 的商。这就是说，$\pi$ 分类这样的 $X_U$ 的\Term{principal-covering}{主覆盖} $Z$：$g_U^{-1}(Z)$ 可延拓为 $S$ 的\Term{etale-covering}{平展覆盖}，且 $Z$ 在 $X_{\bar s}$ 上诱导的\Term{covering}{覆盖}可由某个以 $L$-群为群的\Term{principal-covering}{主覆盖}作结构群\Term{group-extension}{扩张}得到。
\SourcePage{329}
要证明 $\varphi$ 为\Term{isomorphism}{同构}，只需证明这样的\Term{covering}{覆盖} $Z$ 可延拓到整个 $X$。

先证明 $Z$ 可延拓为一个包含 $X_U$ 与 $g(S)$ 的开子集的\Term{etale-covering}{平展覆盖}。设 $W$ 是 $X$ 上的一个\Term{etale}{平展}\Term{scheme}{概形}，其像包含 $X_U$ 与 $g(S)$，置 $W_U=W\times_S U$。由于 $W\to S$ 为 \Term{0-acyclic}{$0$-无环}，且 $\operatorname{prof\text{ét}}_T S\geq2$，有
\[\operatorname{prof\text{ét}}_{W-W_U}W\geq2\]
（SGA 2 XIV 1.13）。因此，若 $Z|_{W_U}$ 可延拓为 $W$ 的\Term{etale-covering}{平展覆盖}，这一延拓在相差唯一\Term{isomorphism}{同构}的意义下唯一。故将 $Z$ 延拓到 $g(S)\cup X_U$ 的某个邻域的问题，在 $g(T)$ 各点附近对\Term{etale-topology}{平展拓扑}具有局部性。若 $t\in T$，置 $x=g(t)$，以 $\bar X$（分别地 $\bar S$）记 $X$ 在 $\bar x$ 处（分别地 $S$ 在 $\bar t$ 处）的\Term{strict-localization}{严格局部化}，$\bar U=U\times_S\bar S$、$\bar X_U=\bar X\times_X X_U$，$\bar g:\bar S\to\bar X$ 为由 $g$ 导出的\Term{morphism}{态射}。只需证明，对 $T$ 的每个点 $t$，$Z$ 在 $\bar X_U$ 上的\Term{inverse-image}{逆像} $\bar Z$ 可延拓到 $\bar X$；等价地，它是\Term{trivial-bundle}{平凡}的。

按 $Z$ 的定义，$\bar Z$ 在 $\bar U$ 上的\Term{inverse-image}{逆像}\Term{trivial-bundle}{平凡}。要证明 $\bar Z$ \Term{trivial-bundle}{平凡}，只需证明它形如 $\bar f_U^*E$，其中 $E$ 为 $\bar U$ 的\Term{principal-covering}{主覆盖}；因为此时
\[E\simeq\bar g_U^*\bar f_U^*E\simeq\bar g_U^*\bar Z,\]
而 $\bar g_U^*\bar Z$ \Term{trivial-bundle}{平凡}，便得所需结论。\Term{morphism}{态射} $\bar f_U$ 为 \Term{0-acyclic}{$0$-无环}且\Term{locally-0-acyclic}{局部 $0$-无环}，因此要证明 $\bar Z$ 来自 $\bar U$，只需证明对位于 $\bar U$ 某点上方并在其上代数的每个\Term{geometric-point}{几何点}——可设为 $\bar s$——$\bar Z|_{X_{\bar s}}$ \Term{trivial-bundle}{平凡}（SGA 4 XV 1.15）。由于 $\bar Z|_{X_{\bar s}}$ 来自一个以 $L$-群为群的\Term{principal-covering}{主覆盖}的结构群\Term{group-extension}{扩张}，这由 $\bar f$ 对 $L$ 为 $1$-无球得到。

于是已经证明，存在 $g(S)\cup X_U$ 的\Term{open-neighborhood}{开邻域} $V$，使 $Z$ 延拓为 $V$ 的\Term{etale-covering}{平展覆盖} $Z_V$。再证明 $Z_V$ 可延拓到整个 $X$。只需说明，对 $X-V$ 的每个点 $x$，有
\[\operatorname{prof\,hop}_xX\geq3.\]
这由假设 e) 得到，因为 $X-V$ 中的点 $x$ 不可能是其所在\Term{fiber}{纤维} $X_t$ 的极大点：$X_t$ 的每个\Term{irreducible-component}{不可约分支}都包含 $g(t)$，故 $X_t$ 的每个极大点都属于 $V$。
\end{proof}

\begin{SGAresult}{推论}{4.8}
沿用 4.7.4 的假设，再设 $\pi_1(S,a)=1$。则有\Term{isomorphism}{同构}
\[\pi_1^L(X,a)\xrightarrow{\sim}\pi_1^L(X_{\bar s},a)_K.\tag{*}\]
特别地，若 $\pi_1^L(X_{\bar s},a)$ 在拓扑意义下\Term{finite-presentation}{有限表示}，且 $K$ 在其上的作用通过一个\Term{finitely-generated}{有限生成}群，则 $\pi_1^L(X,a)$ 在拓扑意义下\Term{finite-presentation}{有限表示}。
\end{SGAresult}
\SourcePage{330}
\Term{isomorphism}{同构} (*) 已在 4.7.4 中证明。设 $\pi_1^L(X_{\bar s},a)$ 为 $n$ 个生成元的自由投-$L$-群 $L(x_1,\ldots,x_n)$ 对元素 $y_1,\ldots,y_p$ 生成的闭\Term{normal-subgroup}{正规子群}的商；并设 $K$ 的作用通过由 $k_1,\ldots,k_q$ 生成的群。对每个 $i\in[1,n]$、$j\in[1,q]$，若以 $z_{ij}$ 记 $L(x_1,\ldots,x_n)$ 中提升元素 $(k_j\cdot x_i)x_i^{-1}$ 的一个元素，则 $\pi_1^L(X_{\bar s},a)_K$ 是 $L(x_1,\ldots,x_n)$ 对元素
\[(y_i)_{i\in[1,p]},\qquad (z_{ij})_{i\in[1,n],\,j\in[1,q]}\]
生成的闭\Term{normal-subgroup}{正规子群}的商。

\begin{SGAresult}{注记}{4.9}
a) 下列情形满足 4.7 的条件 a)—e)：$S$ 是\Term{normal}{正规}\Term{connected}{连通}\Term{scheme}{概形}，$U$ 是 $S$ 的稠密\Term{retrocompact}{相对拟紧}开子集；$f$ 为\Term{proper}{固有}\Term{finite-presentation}{有限表示}\Term{morphism}{态射}，\Term{geometric-fibre}{几何纤维}\Term{connected}{连通}且在 $T$ 的每个点 $t$ 处\Term{irreducible}{不可约}；此外 $f$ \Term{separable}{可分}且在 $X_U\cup g(T)$ 的各点\Term{smooth}{光滑}；$L$ 为不同于 $S$ 的\Term{residue-characteristic}{剩余特征}的素数组成的集合；$X$ 在 $X_t$ 的每个点处\Term{regular}{正则}。事实上，条件 a) 由 SGA 4 XV 4.1、1.4 得到；条件 b)、d) 由 1.4 及 SGA 4 XV 2.1、XVI 5.2 得到；最后，e) 由 SGA XIV 1.11 得到。

b) 推论 4.5 可用于计算特征 $p$ 的\Term{separably-closed}{可分闭}域 $k$ 上\Term{proper}{固有}\Term{smooth}{光滑}曲面 $X$ 的\Term{fundamental-group}{基本群} $\pi_1^{p'}(X)$，其中 $p'$ 为不同于 $p$ 的素数组成的集合。这一方法由 J. P. Murre 告知我们：通过对 $X$ 作\Term{blow-up}{爆破}，归约到存在\Term{fiber}{纤维}化 $X\to\mathbf P^1_k$ 及 $\mathbf P^1_k$ 的一个开子集 $U$，使满足 4.7 假设的情形；详情见 SGA 7。更一般地，同样的方法可用于证明（见上述文献）：若 $X$ 为\Term{connected}{连通}\Term{finite-type}{有限型} $k$-\Term{scheme}{概形}，且 $k$ 的一个\Term{algebraic-closure}{代数闭包}上维数 $\leq\dim X$ 的\Term{finite-type}{有限型}\Term{scheme}{概形}均可强消解奇点（SGA 5 I 3.1.5），则 $\pi_1^{p'}(X)$ 在拓扑意义下\Term{finite-presentation}{有限表示}。
\end{SGAresult}

\SGAsection{5}{附录 I：Abhyankar 引理的各种形式}
本附录给出 Abhyankar 引理的若干变式。

\begin{SGAresult}{命题}{5.1}
设 $X=\operatorname{Spec}A$ 为\Term{regular}{正则}局部\Term{scheme}{概形}，$D=\sum_{1\leq i\leq r}\operatorname{div}f_i$ 为\Term{normal-crossings-divisor}{正规交叉除子}，其中 $f_i$ 是 $A$ 的\Term{maximal-ideal}{极大理想}中的元素，且构成某组\Term{regular-parameters}{正则参数}的一部分。设 $n_i$（$1\leq i\leq r$）为 $\geq0$ 的整数，置
\[X'=X[T_1,\ldots,T_r]/(T_1^{n_1}-f_1,\ldots,T_r^{n_r}-f_r),\qquad U'=U\times_X X'.\]
则 $X'$ \Term{regular}{正则}，$U'$ 为 $X'$ 中\Term{normal-crossings-divisor}{正规交叉除子} $\sum_{1\leq i\leq r}\operatorname{div}T_i$ 的补集。若整数 $n_i$ 与 $X$ 的\Term{residue-characteristic}{剩余特征} $p$ 互素，则 $U'$ 是 $U$ 的\Term{connected}{连通}\Term{etale-covering}{平展覆盖}，且相对于 $D$ \Term{tame-ramification}{温和分歧}（2.3 c)）。
\end{SGAresult}
\SourcePage{331}
\begin{proof}
事实上，$X'$ 是某个\Term{local-ring}{局部环} $A'$ 的谱，其\Term{maximal-ideal}{极大理想}由 $T_1,\ldots,T_r$ 生成；可设 $r=\dim A$。由于 $A'$ 在 $A$ 上\Term{finite}{有限}且\Term{flat}{平坦}，从而维数为 $r$，$A'$ \Term{regular}{正则}（EGA $0_{\mathrm{IV}}$ 17.1.1），且 $T_i$ 构成 $A'$ 的一组\Term{regular-parameters}{正则参数}。设 $n_i$ 与 $p$ 互素。由于各 $f_i$ 在 $U$ 上可逆，$U'$ 在 $U$ 上\Term{etale}{平展}由 I 7.4 得到。此外，$U'$ 相对于 $D$ \Term{tame-ramification}{温和分歧}：设 $x_i$ 为 $V(f_i)$ 的\Term{generic-point}{泛点}，$\bar R$ 为 $R=\mathcal O_{X,x_i}$ 的\Term{strict-localization}{严格局部化}，$\bar K$ 为 $\bar R$ 的\Term{field-of-fractions}{分式域}。表示 $U'|_{\bar K}$ 的 $\bar K$-代数由域 $\bar K[T_i]/(T_i^{n_i}-f_i)$ 经一次\Term{unramified}{非分歧}\Term{group-extension}{扩张}得到，故在 $\bar R$ 上\Term{tame-ramification}{温和分歧}。
\end{proof}

\begin{SGAresult}{命题（绝对 Abhyankar 引理）}{5.2}
设 $X$ 为\Term{regular}{正则}局部\Term{scheme}{概形}，
\[D=\sum_{1\leq i\leq r}\operatorname{div}f_i\]
为 5.1 中那样的\Term{normal-crossings-divisor}{正规交叉除子}，$Y=\operatorname{Supp}D$、$U=X-Y$。设 $V$ 是 $U$ 的一个相对于 $D$ \Term{tame-ramification}{温和分歧}的\Term{etale-covering}{平展覆盖}。若 $x_i$ 是闭集 $V(f_i)$ 的\Term{generic-point}{泛点}，则 $\mathcal O_{X,x_i}$ 为\Term{field-of-fractions}{分式域}为 $K_i$ 的\Term{discrete-valuation-ring}{离散赋值环}，且
\[V|_{K_i}=\operatorname{Spec}\left(\prod_{j\in J_i}L_j\right),\]
其中 $L_j$ 是 $K_i$ 的\Term{finite}{有限}\Term{separable-extension}{可分扩张}。以 $n_j$ 记由 $L_j$ 生成的 Galois \Term{group-extension}{扩张}的\Term{inertia-group}{惯性群}的阶，$n_i$ 为 $j$ 遍历 $J_i$ 时各 $n_j$ 的最小公倍数。置
\[X'=X[T_1,\ldots,T_r]/(T_1^{n_1}-f_1,\ldots,T_r^{n_r}-f_r),\]
$U'=U_{(X')}$、$V'=V_{(X')}$ 等。则 $U'$ 的\Term{etale-covering}{平展覆盖} $V'$ 可在相差唯一\Term{isomorphism}{同构}的意义下唯一延拓为 $X'$ 的\Term{etale-covering}{平展覆盖}，且各 $n_i$ 与 $X$ 的\Term{residue-characteristic}{剩余特征} $p$ 互素。
\end{SGAresult}
\begin{proof}
唯一性来自 $X'$ \Term{normal}{正规}（5.1）：事实上，$X'$ 上延拓 $V'$ 的\Term{etale-covering}{平展覆盖}\Term{isomorphism}{同构}于 $X'$ 在 $V'$ 于 $X'$ \Term{generic-point}{泛点}处的\Term{fiber}{纤维}中的\Term{normalization}{正规化}（10.2）。若 $\bar x'$ 是 $Y'$ 的\Term{geometric-point}{几何点}，以 $\overline{X'}$ 记 $X'$ 在 $\bar x'$ 处的\Term{strict-localization}{严格局部化}，$\overline{V'}=V'_{(\overline{X'})}$，其余记号类似。结合唯一性，利用\Term{descent}{下降}，只需证明对 $Y'$ 的每个\Term{geometric-point}{几何点} $\bar x'$，$\overline{U'}$ 的\Term{etale-covering}{平展覆盖} $\overline{V'}$ 可延拓到 $\overline{X'}$。\Term{regular-scheme}{正则概形} $X'$ 中一个包含所有满足 $\dim\mathcal O_{X',x'}\leq1$ 的点 $x'$ 的开子集的\Term{etale-covering}{平展覆盖}，可延拓到整个 $X$（SGA 2 XIV 1.11），因此甚至只需考虑投影到 $Y'$ 的极大点的点 $\bar x'$。在这样的点，$\overline{V'}$ 可延拓为 $\overline{X'}$ 的\Term{etale-covering}{平展覆盖}由 3.6 得到。

来证明 $n_i$ 与 $p$ 互素。否则例如有 $p\mid n_1$。以
\[X[T_1,\ldots,T_r]/(T_1^{n_1/p}-f_1,T_2^{n_2}-f_2,\ldots,T_r^{n_r}-f_r)\]
替换 $X$ 后，归约到 $X'=X[T_1]/(T_1^p-f_1)$ 的情形。只需证明 $V$ 可延拓为 $X$ 的\Term{etale-covering}{平展覆盖}，因为这样就有 $n_1=1$，与假设矛盾。
\SourcePage{332}
为此可设 $X$ \Term{strictly-local}{严格局部}。设 $Z$ 为 $X$ 中由 $p=0$ 定义的\Term{closed-subscheme}{闭子概形}，$Z_1=Z\cap X_{f_1}$；$Z_1$ 是 $Z$ 的非空开子集。由前面的结论，$U'$ 的\Term{etale-covering}{平展覆盖} $V'$ 可延拓为 $X'$ 的\Term{etale-covering}{平展覆盖} $W'$。设 $W''_1,W''_2$ 分别为 $W'$ 沿两个投影 $X''=X'\times_X X'\rightrightarrows X'$ 的\Term{inverse-image}{逆像}。来证明\Term{descent}{下降}\Term{isomorphism}{同构}
\[u:W''_1|_{U''}\longrightarrow W''_2|_{U''}\]
可延拓为 $X''$-\Term{morphism}{态射} $W''_1\to W''_2$；它必为 $W'$ 相对于 $X'\to X$ 的一个\Term{descent-datum}{下降资料}。设 $Z''$（分别地 $Z''_1$）为 $Z$（分别地 $Z_1$）在 $X''$ 中的\Term{inverse-image}{逆像}。由于 $Z''\to Z$ 为\Term{radicial}{根性}\Term{morphism}{态射}，存在延拓 $u|_{Z''_1}$ 的\Term{isomorphism}{同构} $v:W''_1|_{Z''}\to W''_2|_{Z''}$。又由于 $X$ 为 Hensel \Term{scheme}{概形}，有\Term{bijection}{双射}
\[\operatorname{Hom}_{X''}(W''_1,W''_2)\simeq\operatorname{Hom}_{Z''}(W''_1|_{Z''},W''_2|_{Z''}),\]
从而有提升 $v$ 的\Term{morphism}{态射} $w:W''_1\to W''_2$。$X''$ 中使 $u,w$ 在其上相等的开\Term{closed-subscheme}{闭子概形}包含 $Z''_1$，因而等于 $X''$；于是 $V$ 可延拓到 $X$。
\end{proof}

\textbf{5.3.0.} 沿用 5.2 的假设和记号，再设 $X$ \Term{strictly-local}{严格局部}。由上述结论，$U$ 的每个相对于 $D$ \Term{tame-ramification}{温和分歧}的\Term{connected}{连通}\Term{etale-covering}{平展覆盖}，都是一个形如
\[U'=U[T_1,\ldots,T_r]/(T_1^{n_1}-f_1,\ldots,T_r^{n_r}-f_r)\]
的\Term{tame-ramification}{温和分歧}\Term{covering}{覆盖}的商，其中 $n_i$ 是与 $p$ 互素的整数。以 $\boldsymbol\mu_n$ 记 $U$ 的 $n$ 次单位根群。$U'$ 的 $U$-\Term{automorphism}{自同构}群正是 $\boldsymbol\mu_{n_1}\times\cdots\times\boldsymbol\mu_{n_r}$；$n_i$ 次单位根 $\xi_i$ 将 $T_i$ 变为 $\xi_iT_i$，由此作用于 $U'$。故有下列结论。

\begin{SGAresult}{推论}{5.3}
设 $X$ 为\Term{residue-characteristic}{剩余特征} $p\geq0$ 的\Term{regular}{正则}\Term{strictly-local}{严格局部}\Term{scheme}{概形}，$D=\sum_{1\leq i\leq n}\operatorname{div}f_i$ 为 $X$ 上的\Term{normal-crossings-divisor}{正规交叉除子}，$U=X-\operatorname{Supp}D$。置
\[\widetilde U=\varprojlim_{(n_i)}U[T_1,\ldots,T_r]/(T_1^{n_1}-f_1,\ldots,T_r^{n_r}-f_r),\]
\Term{inverse-limit}{投射极限}按由与 $p$ 互素的正整数 $n_i$ 的族组成的、以整除关系排序的滤过有序集取得。则 $\widetilde U$ 是 $U$ 的\Term{tame-ramification}{温和分歧}\Term{universal-covering}{万有覆盖}。因此，$U$ 的\Term{tame-fundamental-group}{温和分歧基本群}为
\[\pi_1^t(U)\simeq\prod_{\ell\ne p}\mathbf Z_\ell[1]^r\qquad\text{（典范同构）},\]
其中 $\mathbf Z_\ell[1]=\varprojlim_{n>0}\boldsymbol\mu_{\ell^n}$。群 $\pi_1^t(U)$ 非典范\Term{isomorphism}{同构}于 $\prod_{\ell\ne p}\mathbf Z_\ell^r$。
\end{SGAresult}

\SourcePage{333}
\begin{SGAresult}{命题}{5.4}
设 $f:X\to S$ 为\Term{scheme}{概形}\Term{morphism}{态射}，$D=\sum_{1\leq i\leq r}\operatorname{div}f_i$ 为一个相对于 $S$ 具有\Term{normal-crossings}{正规交叉}的\Term{divisor}{除子}（2.1）。对 $Y=\operatorname{Supp}D$ 的每个点 $x$，若 $I(x)\subset[1,r]$ 是满足 $f_i(x)=0$ 的 $i$ 所成的集合，则\Term{subscheme}{子概形} $V((f_i)_{i\in I(x)})$ 在 $S$ 上\Term{smooth}{光滑}，在 $X$ 中的余维为 $\operatorname{card}I(x)$。置 $U=X-Y$。设 $x\in Y$，$X_1=\operatorname{Spec}\mathcal O_{X,x}$、$U_1=U\times_X X_1$，$n_i$（$i\in I(x)$）为整数，并置
\[X'=X[T_i]_{i\in I(x)}/(T_i^{n_i}-f_i).\]
若 $x'$ 是 $X'$ 中位于 $x$ 上方的点，则 $X'$ 在 $x'$ 处在 $S$ 上\Term{smooth}{光滑}。若整数 $n_i$ 与 $k(x)$ 的特征 $p$ 互素，则 $U'_1=U_1\times_X X'$ 是 $U_1$ 的\Term{connected}{连通}\Term{etale-covering}{平展覆盖}，且在 $X_1$ 上相对于 $S_1$ \Term{tame-ramification}{温和分歧}（2.1.1）。
\end{SGAresult}
\begin{proof}
若 $s=f(x)$，则\Term{geometric-fibre}{几何纤维} $X'_{\bar s}$ 在 $x'$ 处\Term{regular}{正则}（5.1）；由于 $X'$ 在 $Y$ 的邻域内在 $S$ 上\Term{flat}{平坦}，这说明 $X'$ 在 $x'$ 处在 $S$ 上\Term{smooth}{光滑}（EGA IV 12.1.6）。若 $n_i$ 与 $p$ 互素，则 $U'_1$ 是 $U_1$ 的\Term{etale-covering}{平展覆盖}（7.4）。它在 $X_1$ 上相对于 $S$ \Term{tame-ramification}{温和分歧}，因为在 $S$ 的每个点的\Term{geometric-fibre}{几何纤维}上均如此（5.1）。最后，$U'_1$ \Term{connected}{连通}由 SGA 4 XVI 3.2 得到。
\end{proof}

\begin{SGAresult}{命题（相对 Abhyankar 引理）}{5.5}
设 $X$ 是 $S$-\Term{scheme}{概形}，$D$ 为 5.4 中那样的相对于 $S$ 具有\Term{normal-crossings}{正规交叉}的\Term{divisor}{除子}。置 $Y=X-\operatorname{Supp}D$、$U=X-Y$；设 $x\in Y$，$X_1$ 为 $X$ 在 $x$ 上方一个\Term{geometric-point}{几何点}处的\Term{strict-localization}{严格局部化}，$U_1=U\times_X X_1$，$V_1$ 为 $U_1$ 的一个\Term{etale-covering}{平展覆盖}。设对 $S$ 的每个极大点 $s$，$V_{1s}$ 在 $X_{1s}$ 上相对于 $s$ \Term{tame-ramification}{温和分歧}。则存在与 $k(x)$ 的特征 $p$ 互素的整数 $n_i$（$i\in I(x)$），使得置
\[X'_1=X_1[T_i]_{i\in I(x)}/(T_i^{n_i}-f_i),\]
$U'_1=U_1\times_{X_1}X'_1$ 等后，$U'_1$ 的\Term{etale-covering}{平展覆盖} $V'_1$ 可在相差唯一\Term{isomorphism}{同构}的意义下唯一延拓为 $X'_1$ 的\Term{etale-covering}{平展覆盖}。特别地，$V_1$ 在 $X_1$ 上相对于 $S$ \Term{tame-ramification}{温和分歧}。
\end{SGAresult}
\begin{proof}
可设 $S$ 为\Term{closed-point}{闭点}为 $f(x)$ 的\Term{locally-noetherian}{局部 Noether} \Term{scheme}{概形}。对 $S$ 的每个极大点 $s$ 和每个 $i\in I(x)$，以 $x_i$ 记\Term{fiber}{纤维} $X_{1s}$ 中闭集 $V(f_i)$ 的\Term{generic-point}{泛点}。\Term{local-ring}{局部环} $(\mathcal O_{X_1,x_i})_{\mathrm{red}}$ 为\Term{field-of-fractions}{分式域}为 $K_i$ 的\Term{discrete-valuation-ring}{离散赋值环}，并有
\[V_1|_{K_i}=\operatorname{Spec}\left(\prod_{j\in I(x_i)}L_j\right),\]
其中 $L_j$ 为 $K$ 的\Term{finite}{有限}\Term{separable-extension}{可分扩张}。以 $n_j$ 记由 $L_j$ 生成的 Galois \Term{group-extension}{扩张}的\Term{inertia-group}{惯性群}的阶，$n_i$ 为 $s$ 遍历 $S$ 的极大点且 $j\in I(x_i)$ 时各 $n_j$ 的最小公倍数。

按此选择 $n_i$，来证明 $V'_1$ 唯一延拓为 $X'_1$ 的\Term{etale-covering}{平展覆盖}。唯一性来自：$X'$ 在 $Y$ 的各点在 $S$ 上\Term{smooth}{光滑}，故 $\operatorname{prof\text{ét}}_{Y'_1}(X'_1)\geq2$（SGA 4 XVI 3.2 或 SGA 2 XIV 1.19）。设 $x'_1$ 为 $Y'_1$ 的一个点，$\bar x'_1$ 为其上方的\Term{geometric-point}{几何点}，以 $\overline{X'_1}$ 记 $X'_1$ 在 $\bar x'_1$ 处的\Term{strict-localization}{严格局部化}，
\SourcePage{334}
$\overline{U'_1}=U'_{1(\overline{X'_1})}$，其余记号类似。结合唯一性，利用\Term{descent}{下降}，只需证明 $\overline{V'_1}$ 可延拓到 $\overline{X'_1}$。此外，只需取 $Y'_1$ 的极大点 $x'_1$：这样就有 $X'_1$ 的一个包含 $Y'_1$ 所有极大点的开子集 $W'_1$，使 $V'_1$ 延拓到 $W'_1$。若 $Z'_1=X'_1-W'_1$，则当 $s$ 为 $S$ 的极大点时，有 $\operatorname{codim}(Z'_{1s},X'_{1s})\geq2$；当 $s$ 不是极大点时，有 $\operatorname{codim}(Z'_{1s},X'_{1s})\geq1$、$\operatorname{prof\text{ét}}_s(S)\geq1$。于是由 SGA 2 XIV 1.20，$V'_1$ 可延拓到整个 $X'_1$。而在 $Y'_1$ 的一个极大点上方的\Term{geometric-point}{几何点} $\bar x'_1$ 处，$\overline{X'_1}_{\mathrm{red}}$ 是\Term{discrete-valuation-ring}{离散赋值环}的谱，$\overline{V'_1}$ 可延拓到 $\overline{X'_1}$ 便由 X 3.6 得到。

来证明 $n_i$ 与 $p$ 互素。否则存在 $i_0\in I(x)$，使 $p\mid n_{i_0}$。以
\[X_1[T_{i_0},T_i]_{i\in I(x)}/(T_i^{n_{i_0}/p}-f_{i_0},T_i^{n_i}-f_i)\]
替换 $X$ 后，归约到 $X'_1=X_1[T]/T^p-f_{i_0}$ 的情形。由前面的结论，$U'_1$ 的\Term{etale-covering}{平展覆盖} $V'_1$ 延拓为 $X'_1$ 的\Term{etale-covering}{平展覆盖} $E'_1$。设 $\eta$ 为 $S$ 的\Term{closed-point}{闭点}；由于 $X'_{1\eta}\to X_{1\eta}$ 为\Term{radicial}{根性}\Term{morphism}{态射}，$V_{1\eta}$ 延拓为 $X_{1\eta}$ 的\Term{etale-covering}{平展覆盖} $E_{1\eta}$。像 5.2 那样，$E'_1$ 带有相对于 $X'_1\to X_1$ 的\Term{descent-datum}{下降资料}，它延拓了 $E'_1|_{U'_1}$ 上的自然\Term{descent-datum}{下降资料}。故 $V_1$ 可延拓到 $X_1$；但这蕴含 $n_{i_0}=1$，与假设 $n_{i_0}=p$ 矛盾。
\end{proof}

\begin{SGAresult}{推论}{5.6}
设 $X$ 是 $S$-\Term{scheme}{概形}，$D=\sum_{1\leq i\leq r}\operatorname{div}f_i$ 是 5.4 中那样的相对于 $S$ 具有\Term{normal-crossings}{正规交叉}的\Term{divisor}{除子}。设 $\bar x$ 是 $X$ 的\Term{geometric-point}{几何点}，$\bar X$ 为 $X$ 在 $\bar x$ 处的\Term{strict-localization}{严格局部化}，$\bar Y=Y_{(\bar X)}$、$\bar U=\bar X-\bar Y$，并置
\[\widetilde U=\varprojlim_{(n_i)}\bar U[T_i]_{i\in I(x)}/(T_i^{n_i}-f_i),\]
其中\Term{inverse-limit}{投射极限}按与 $k(x)$ 的特征 $p$ 互素的正整数 $n_i$ 的族所成的滤过集取得。则 $\widetilde U$ 是 $\bar U$ 相对于 $S$ 的\Term{tame-ramification}{温和分歧}\Term{universal-covering}{万有覆盖}。因此，$\bar U$ 的\Term{tame-fundamental-group}{温和分歧基本群}为
\[\pi_1^t(\bar U)\simeq\prod_{\ell\ne p}\mathbf Z_\ell[1]^{I(x)}\qquad\text{（典范同构）}.\]
群 $\pi_1^t(\bar U)$ 非典范\Term{isomorphism}{同构}于 $\prod_{\ell\ne p}\mathbf Z_\ell^{I(x)}$。
\end{SGAresult}

\begin{SGAresult}{注记}{5.6.1}
设 $X$ 是 $S$-\Term{scheme}{概形}，$D=\sum_{1\leq i\leq r}\operatorname{div}f_i$ 是 5.4 中那样的相对于 $S$ 具有\Term{normal-crossings}{正规交叉}的\Term{divisor}{除子}，$U=X-\operatorname{Supp}D$。对每个子集 $I\subset[1,r]$，置
\[X_I=\left(\bigcap_{i\in I}V(f_i)\right)\cap\left(\bigcap_{i\in\complement I}X_{f_i}\right).\]
\SourcePage{335}
设 $p$ 为素数或零，$Z$ 是 $X_I$ 的一个子集，其所有点的特征均为 $p$。置
\[\widetilde U_I=\varprojlim_{(n_i)}U[T_i]_{i\in I}/(T_i^{n_i}-f_i),\]
其中\Term{inverse-limit}{投射极限}按与 $p$ 互素的正整数 $n_i$ 的族所成的滤过集取得。则对 $Z$ 的每个\Term{geometric-point}{几何点} $\bar x$，$\widetilde U_I$ 在 $\bar U$ 上的\Term{inverse-image}{逆像}识别为 $\bar U$ 的\Term{tame-ramification}{温和分歧}\Term{universal-covering}{万有覆盖}。
\end{SGAresult}

\begin{SGAresult}{推论}{5.7}
沿用 5.6 的记号。设 $\bar S$ 为 $S$ 在 $\bar x$ 处的\Term{strict-localization}{严格局部化}，
\[\bar g:\bar U\longrightarrow\bar S,\qquad\widetilde g:\widetilde U\longrightarrow\bar S\]
为典范\Term{morphism}{态射}。则 $\bar g$ 与 $\widetilde g$ 均为 \Term{0-acyclic}{$0$-无环}（SGA 4 XV 1.3）。设 $G$ 为 $\bar S$ 上的\Term{constructible}{可构造}\Term{group-sheaf}{群层}，$F=\bar g^*G$，$P$ 为 $F$-\Term{torsor}{挠子}。则 $P$ 在 $\bar X$ 上相对于 $\bar S$ \Term{tame-ramification}{温和分歧}，当且仅当它在 $\widetilde U$ 上的\Term{inverse-image}{逆像} $\widetilde P$ \Term{trivial-bundle}{平凡}。
\end{SGAresult}
\begin{proof}
事实上，对每个 $\overline{X'}=\bar X[T_i]_{i\in I(x)}/(T_i^{n_i}-f_i)$，其中 $n_i>0$ 且与 $p$ 互素，\Term{morphism}{态射} $\overline{f'}:\overline{X'}\to\bar S$ 为 \Term{0-acyclic}{$0$-无环}。因此，$\overline{f'}$ 在 $\bar S$ 各点处的\Term{geometric-fibre}{几何纤维}\Term{connected}{连通}，甚至\Term{irreducible}{不可约}。故 $\overline{g'}:\overline{U'}\to\bar S$ 的\Term{geometric-fibre}{几何纤维}亦如此，说明 $\overline{g'}$，从而 $\widetilde g$，为 \Term{0-acyclic}{$0$-无环}（SGA 4 XV 1.16）。

显然，$\bar U$ 上一个在 $\widetilde U$ 上的\Term{inverse-image}{逆像}\Term{trivial-bundle}{平凡}的 $F$-\Term{torsor}{挠子} $P$，在 $\bar X$ 上相对于 $\bar S$ \Term{tame-ramification}{温和分歧}。反之，设 $P$ 在 $\bar X$ 上相对于 $\bar S$ \Term{tame-ramification}{温和分歧}，来证明其在 $\widetilde U$ 上的\Term{inverse-image}{逆像}\Term{trivial-bundle}{平凡}。

由 SGA 4 IX 2.14(ii)，可取\Term{finite-morphism}{有限态射} $n:S_1\to\bar S$、$S_1$ 上的常值\Term{group-sheaf}{群层} $C$ 和\Term{monomorphism}{单态射} $G\to n_*C$。考虑下列由\Term{cartesian-square}{笛卡尔方块}组成的\Term{commutative-diagram}{交换图}：
\[
\begin{tikzcd}
\widetilde U_1\arrow[r]\arrow[d,"r"']&U_1\arrow[r,"g_1"]\arrow[d,"q"]&S_1\arrow[d,"n"]\\
\widetilde U\arrow[r]&\bar U\arrow[r,"\bar g"']&\bar S.
\end{tikzcd}
\]
设 $C_1$（分别地 $\widetilde C_1$）为 $C$ 在 $U_1$（分别地 $\widetilde U_1$）上的\Term{inverse-image}{逆像}。有\Term{commutative-diagram}{交换图}，其中 $i,j$ 为\Term{isomorphism}{同构}（SGA 4 VIII 5.8）：
\[
\begin{tikzcd}[column sep=large]
H^1(\bar U,q_*C_1)\arrow[r,"i"]\arrow[d]&H^1(U_1,C_1)\arrow[d]\\
H^1(\widetilde U,r_*\widetilde C_1)\arrow[r,"j"']&H^1(\widetilde U_1,\widetilde C_1).
\end{tikzcd}\tag{*}
\]
设 $Q$ 为 $P$ 沿结构群\Term{group-extension}{扩张} $F\to q_*C_1$ 得到的 $q_*C_1$-\Term{torsor}{挠子}。由 2.1.4，$Q$ 在 $\bar X$ 上相对于 $\bar S$ \Term{tame-ramification}{温和分歧}。
\SourcePage{336}
通过 $i$，\Term{torsor}{挠子} $Q$ 对应于 $C_1$-\Term{torsor}{挠子} $Q_1$；显然，$Q_1$ 在 $X_1=\bar X\times_{\bar S}S_1$ 上相对于 $S_1$ \Term{tame-ramification}{温和分歧}。由 5.6，$Q_1$ 在 $\widetilde U_1$ 上的\Term{inverse-image}{逆像} $\widetilde Q_1$ \Term{trivial-bundle}{平凡}；图 (*) 便说明 $Q$ 在 $\widetilde U$ 上的\Term{inverse-image}{逆像} $\widetilde Q$ \Term{trivial-bundle}{平凡}。

考虑下列\Term{commutative-diagram}{交换图}，其第二行\Term{exact}{正合}（SGA 4 XII 3.1）：
\begin{center}
\resizebox{.9\textwidth}{!}{\begin{tikzcd}[ampersand replacement=\&,column sep=large]
H^0(\bar S,n_*C/G)\arrow[r]\arrow[d,"k"']\&H^1(\bar S,G)=1\arrow[d]\&\\
H^0(\widetilde U,r_*\widetilde C_1/\widetilde F)\arrow[r]\&H^1(\widetilde U,\widetilde F)\arrow[r]\&H^1(\widetilde U,r_*\widetilde C_1).
\end{tikzcd}}\quad(**)
\end{center}
由于 $\widetilde U\to\bar S$ 为 \Term{0-acyclic}{$0$-无环}，$k$ 是\Term{isomorphism}{同构}。由 (**) 即得 $\widetilde P$ \Term{trivial-bundle}{平凡}。
\end{proof}

\SGAsection{6}{附录 II：叠的正像的有限性定理}
\begin{SGAresult}{命题}{6.1}
设 $S$ 是\Term{locally-noetherian}{局部 Noether} \Term{scheme}{概形}，$f:X\to S$ 为\Term{morphism}{态射}。若 $S'$ 是 $S$-\Term{scheme}{概形}，以 $X'$（分别地 $f'$ 等）记 $X$（分别地 $f$ 等）沿 $S'\to S$ 的\Term{inverse-image}{逆像}。设对 $S$ 上每个\Term{etale}{平展}\Term{scheme}{概形} $S'$，$X'$ 上每个\Term{constructible}{可构造}\Term{set-sheaf}{集合层} $F$ 的 $f'_*F$ \Term{constructible}{可构造}，且 $X'$ 上每个\Term{constructible}{可构造}\Term{group-sheaf}{群层} $F$ 的 $R^1f'_*F$ \Term{constructible}{可构造}。若 $\Phi$ 是 $X$ 上的 $1$-\Term{constructible}{可构造}叠（0），则 $f_*\Phi$ 是 $1$-\Term{constructible}{可构造}的。
\end{SGAresult}
\begin{proof}
对 $S$ 上每个\Term{etale}{平展}\Term{scheme}{概形} $S'$ 和 $(f_*\Phi)_{S'}$ 中每个对象 $x$，有\Term{isomorphism}{同构}
\[\underline{\operatorname{Aut}}_{S'}(x)\simeq f'_*\underline{\operatorname{Aut}}_{X'}(x),\]
其中等式右侧把 $x$ 视为 $\Phi_{X'}$ 中的对象。故由假设，$f_*\Phi$ \Term{constructible}{可构造}。以 $S\Phi$ 记 $\Phi$ 的\Term{maximal-subgerbe}{极大子凝集}层 [1, III 2.1.7]。由于 $f_*(S\Phi)$ \Term{constructible}{可构造}，可对它应用 SGA 4 IX 2.7，而 $f_*\Phi$ 的 $1$-\Term{constructible}{可构造}性由下列引理得到。
\end{proof}

\begin{SGAresult}{引理}{6.1.1}
设 $S$ 是\Term{locally-noetherian}{局部 Noether} \Term{scheme}{概形}，$f:X\to S$ 为\Term{morphism}{态射}，$\Phi$ 为 $X$ 上的叠。设给定 $S$ 上一个由\Term{finite-type}{有限型}\Term{etale}{平展} $S$-\Term{scheme}{概形} $T$ 表示的层、\Term{surjectivity}{满射}\Term{morphism}{态射}
\[a:T\longrightarrow f_*(S\Phi),\]
以及\Term{fiber}{纤维} $\Phi_{X_T}$ 中的一个对象 $p$（其中 $X_T=X\times_S T$），使它在 $f_*(S\Phi)(T)=S\Phi(X_T)$ 中定义的元素等于 $T(T)$ 的恒等\Term{section}{截面}经 $a$ 所得的像 $q$。设 $f_T:X_T\to T$ 为典范\Term{morphism}{态射}，并设层 $R^1f_{T*}(\underline{\operatorname{Aut}}_{X_T}(p))$ \Term{constructible}{可构造}，则 $S(f_*\Phi)$ 也\Term{constructible}{可构造}。
\end{SGAresult}
\SourcePage{337}
\begin{proof}
典范\Term{morphism}{态射} $f^*f_*\Phi\to\Phi$ 给出\Term{morphism}{态射}
\[S(f^*f_*\Phi)\simeq f^*(S(f_*\Phi))\longrightarrow S\Phi,\]
进而给出典范\Term{morphism}{态射}
\[\varphi:S(f_*\Phi)\longrightarrow f_*(S\Phi).\]
设 $F=S(f_*\Phi)$，$G$ 为 $F$ 经 $\varphi$ 的像。由 SGA 4 IX 2.9，$G$ 是\Term{constructible}{可构造}层。

只需证明：对 $S$ 的每个点 $s$，$\overline{\{s\}}$ 中存在非空开子集 $U$，使 $F|_U$ \Term{locally-constant}{局部常值}且\Term{constructible}{可构造}。设 $s\in S$，$\bar s$ 为其上方的\Term{geometric-point}{几何点}，$\bar q_1,\ldots,\bar q_n$ 为 $G_{\bar s}$ 的各元素。由 $T$ 的定义，存在 $S$-\Term{morphism}{态射} $h_i:\bar s\to S'$，使 $\bar q_i=h_i^*(q)$。设 $S'$ 为 $n$ 个\Term{isomorphism}{同构}于 $T$ 的\Term{scheme}{概形}在 $S$ 上的\Term{fiber-product}{纤维积}，$\bar s\to S'$ 为 $h_i$ 的\Term{fiber-product}{纤维积}，$X'=X\times_S S'$，$q_i$（分别地 $p_i$）为 $q$（分别地 $p$）沿 $S'\to T$ 的第 $i$ 个投影的\Term{inverse-image}{逆像}。若 $\Psi_i$ 为 $p_i$ 生成的 $\Phi|_{X'}$ 的\Term{maximal-subgerbe}{极大子凝集}，则层 $F_i=R^1f'_*\underline{\operatorname{Aut}}_{X'}(p_i)$ 正是 $f'_*\Psi_i$ 的\Term{maximal-subgerbe}{极大子凝集}层 $S(f'_*\Psi_i)$。特别地，典范\Term{injectivity}{单射} $\Psi_i\to\Phi|_{X'}$ 给出\Term{morphism}{态射}
\[\alpha_i:F_i\longrightarrow F|_{S'}.\]
来证明 $\alpha_i$ 是从 $F_i$ 到 $F|_{S'}$ 中 $q_i$ 的\Term{inverse-image}{逆像}的\Term{bijection}{双射}。对 $S'$ 上每个\Term{etale}{平展}\Term{scheme}{概形} $S''$，$F_i(S'')$ 的每个\Term{section}{截面} $y$ 在 $F(S'')$ 中的像为 $q_i|_{S''}$，因为在 $S''$ 的\Term{etale-topology}{平展拓扑}下，$y$ 局部由与 $p_i|_{X''}$ \Term{isomorphism}{同构}的 $\Phi_{X''}$ 中的对象 $x$ 定义。反之，若 $y\in F(S'')$ 在 $F(S'')$ 中的像为 $q_i|_{S''}$，则在 $S''$ 的\Term{etale-topology}{平展拓扑}下，$y$ 局部由\Term{isomorphism}{同构}于 $p_i$ 的 $\Phi_{X''}$ 中的对象 $x$ 定义；于是 $x$ 是 $\Psi_{iX''}$ 的对象，因而 $y\in F_i(S'')$。

利用下文 6.1.2 即可完成证明。事实上，可取 $\overline{\{s\}}$ 的一个\Term{open-neighborhood}{开邻域} $U'$，使 $q_1|_{U'},\ldots,q_n|_{U'}$ 为 $G(U')$ 的\Term{section}{截面}并生成该层。由于 $F_i|_{U'}$ 与 $G|_{U'}$ \Term{constructible}{可构造}，由 6.1.2，$F|_{U'}$ 亦\Term{constructible}{可构造}。再缩小 $U$，可使 $F|_U$ \Term{locally-constant}{局部常值}，证明完成。
\end{proof}

\begin{SGAresult}{引理}{6.1.2}
设 $S$ 是\Term{locally-noetherian}{局部 Noether} \Term{scheme}{概形}，$F\to G$ 是 $S$ 上\Term{group-sheaf}{群层}的一个\Term{surjectivity}{满射}\Term{morphism}{态射}。设 $q_i$ 是 $G$ 在 $X$ 上的一族\Term{finite}{有限}个\Term{section}{截面}，它们生成 $G$；对每个 $i$，设 $F_i$ 为 $F$ 中 $q_i$ 的\Term{inverse-image}{逆像}\Term{subsheaf}{子层}。若 $G$ 与各 $F_i$ \Term{constructible}{可构造}，则 $F$ 也\Term{constructible}{可构造}。
\end{SGAresult}
\begin{proof}
要证明 $F$ \Term{constructible}{可构造}，只需证明：对 $S$ 的每个点 $s$，$\overline{\{s\}}$ 中存在\Term{open-neighborhood}{开邻域} $U$，使 $F|_U$ \Term{locally-constant}{局部常值}且\Term{constructible}{可构造}。取 $S$ 的一点 $s$。由于层 $F_i,G$ \Term{constructible}{可构造}，可取 $\overline{\{s\}}$ 中的\Term{open-neighborhood}{开邻域} $U$，使 $F_i|_U$ 与 $G|_U$ 都\Term{locally-constant}{局部常值}。来证明 $F|_U$ \Term{locally-constant}{局部常值}。由 SGA 4 IX 2.13(i)，
\SourcePage{338}
只需证明：若 $\bar s$ 是位于 $s$ 上方的\Term{geometric-point}{几何点}，$\widetilde s$ 是 $U$ 的\Term{geometric-point}{几何点}，$\bar s\to\widetilde s$ 是\Term{specialization-morphism}{特殊化态射}，则典范\Term{morphism}{态射}
\[F_{\widetilde s}\longrightarrow F_{\bar s}\]
为\Term{bijection}{双射}。考虑\Term{commutative-diagram}{交换图}
\[
\begin{tikzcd}
(F_i)_{\widetilde s}\arrow[r,"\widetilde a_i"]\arrow[d,"\wr"']&F_{\widetilde s}\arrow[r,"\widetilde a"]\arrow[d,"b"]&G_{\widetilde s}\arrow[d,"\wr"]\\
(F_i)_{\bar s}\arrow[r,"\bar a_i"']&F_{\bar s}\arrow[r,"\bar a"']&G_{\bar s}.
\end{tikzcd}
\]
设 $\bar q_i$（分别地 $\widetilde q_i$）为 $q_i$ 在 $G_{\bar s}$（分别地 $G_{\widetilde s}$）中的\Term{inverse-image}{逆像}。\Term{morphism}{态射} $\bar a,\widetilde a$ \Term{surjectivity}{满射}，$\bar a_i$（分别地 $\widetilde a_i$）诱导从 $(F_i)_{\bar s}$（分别地 $(F_i)_{\widetilde s}$）到 $\bar a^{-1}(q_i)$（分别地 $\widetilde a^{-1}(q_i)$）的\Term{bijection}{双射}。由上图可知 $b$ 为\Term{isomorphism}{同构}。
\end{proof}

\begin{SGAresult}{推论}{6.2}
设 $S$ 是\Term{locally-noetherian}{局部 Noether} \Term{scheme}{概形}，$f:X\to S$ 为\Term{proper-morphism}{固有态射}。若 $\Phi$ 是 $X$ 上的 $1$-\Term{constructible}{可构造}叠，则 $f_*\Phi$ 为 $1$-\Term{constructible}{可构造}叠。
\end{SGAresult}
结合 2.4 2)，6.1 的证明也给出下列结果。

\begin{SGAresult}{推论}{6.3}
设 $S$ 是\Term{locally-noetherian}{局部 Noether} \Term{scheme}{概形}，$f:X\to S$ 为\Term{morphism}{态射}，$D$ 为 $X$ 上一个相对于 $S$ 具有\Term{normal-crossings}{正规交叉}的\Term{divisor}{除子}（2.1），$Y=\operatorname{Supp}D$、$U=X-Y$，$i:U\to X$ 为典范\Term{immersion}{浸入}。设 $\Phi$ 是 $U$ 上的一个叠，在 $X$ 与 $S$ 的\Term{etale-topology}{平展拓扑}下，局部为 $S$ 上某个 $1$-\Term{constructible}{可构造}叠 $\Psi$ 的\Term{inverse-image}{逆像}。则叠 $i_*^t\Phi$ 是 $1$-\Term{constructible}{可构造}的。
\end{SGAresult}

\SGAsection{7}{参考文献}
\begin{enumerate}[label={[\arabic*]},leftmargin=2em]
\item J. Giraud，\textit{Cohomologie non abélienne de degré 2}（《二次非交换\Term{cohomology}{上同调}》），博士论文，巴黎，1966。
\item J. Giraud，\textit{Cohomologie non abélienne}（《非交换\Term{cohomology}{上同调}》），Springer，柏林—纽约，1971。
\item H. Seifert、W. Threlfall，\textit{Lehrbuch der Topologie}（《拓扑学教程》），Chelsea，纽约，1934。
\item J.-P. Serre，\textit{Cohomologie galoisienne}（《Galois \Term{cohomology}{上同调}》），Lecture Notes 第 5 卷，Springer，柏林，1964。
\item J.-P. Serre，\textit{Corps locaux}（《局部域》），Hermann，巴黎，1962。
\end{enumerate}
